\documentclass[11pt,twocolumn,a4paper]{article}
\usepackage{fontspec}
\setmainfont{Noto Serif CJK SC}[BoldFont={Noto Serif CJK SC Bold}]
\setsansfont{Noto Sans CJK SC}
\setmonofont{DejaVu Sans Mono}[Scale=.85]
\XeTeXlinebreaklocale "zh"
\XeTeXlinebreakskip=0pt plus .8pt
\usepackage{amsmath,amssymb,graphicx,geometry,booktabs,array,tabularx,multirow,longtable}
\usepackage{caption,subcaption,float,dblfloatfix,placeins}
\usepackage[numbers,square,sort&compress]{natbib}
\renewcommand{\bibfont}{\fontspec{Latin Modern Roman}\fontsize{9}{10.6}\selectfont}
\usepackage[unicode,colorlinks=true,allcolors=black]{hyperref}
\usepackage{titlesec,fancyhdr,balance,xcolor}
\geometry{left=19mm,right=19mm,top=20mm,bottom=21mm,columnsep=7mm}
\setlength{\parindent}{1em}\setlength{\parskip}{2pt plus .4pt}
\setlength{\textfloatsep}{10pt plus2pt minus2pt}\setlength{\dbltextfloatsep}{11pt plus2pt minus2pt}
\setlength{\abovedisplayskip}{6pt plus1pt minus1pt}\setlength{\belowdisplayskip}{6pt plus1pt minus1pt}
\renewcommand{\baselinestretch}{1.10}
\renewcommand{\topfraction}{.92}\renewcommand{\dbltopfraction}{.98}\renewcommand{\textfraction}{.05}
\renewcommand{\floatpagefraction}{.8}\setcounter{topnumber}{2}\setcounter{dbltopnumber}{1}
\captionsetup{font=small,labelfont=bf,labelsep=quad,skip=4pt,justification=raggedright,singlelinecheck=false}
\renewcommand{\figurename}{图}\renewcommand{\tablename}{表}\renewcommand{\refname}{参考文献}
\titleformat{\section}{\large\bfseries}{\thesection}{.55em}{}
\titleformat{\subsection}{\normalsize\bfseries}{\thesubsection}{.5em}{}
\titlespacing*{\section}{0pt}{9pt plus2pt minus1pt}{4pt}\titlespacing*{\subsection}{0pt}{7pt plus2pt minus1pt}{3pt}
\pagestyle{fancy}\fancyhf{}\fancyfoot[C]{\small\thepage}\renewcommand{\headrulewidth}{0pt}
\setlength{\emergencystretch}{2em}
\providecommand{\tightlist}{\setlength{\itemsep}{0pt}\setlength{\parskip}{0pt}}
\newcommand{\real}[1]{#1}
\begin{document}\raggedbottom
\twocolumn[\begin{center}
{\fontsize{17}{20}\selectfont\bfseries Lost in Contrast: Preserving Abstention\\in Multimodal Contrastive Decoding\par}
\vspace{7pt}{\fontsize{12}{16}\selectfont 保留弃权指令的多模态对比解码\par}\vspace{10pt}
\end{center}]

\section*{摘要}
视觉对比解码增强答案对图像证据的依赖，却会把已有依据的弃权改写为错误答案。本文揭示这一冲突来自两个连续操作：正常分布先决定弃权词元能否参与竞争，参考分布再通过对指令的响应改变其相对优势。即使两条分支收到相同的弃权要求，参考侧更强的响应也会在相减时抵消正常侧的弃权倾向。我们据此提出指令保持的三条件解码，将任务引导与视觉差分分别构造，使对比校正在纠正具体答案的同时保留正常视觉条件中的指令作用。该构造得到扰动图像参考的 IP-VCD 和去图像参考的 IP-M3ID。在 Food-101 上，IP-VCD 提高九个模型中八个的联合决策效用；在包含真实不可回答问题的 VizWiz 上，两种实例均在九模型中获益，平均提升分别为13.72和14.66个百分点。四路分布、候选次序翻转与完整回复的对应表明，视觉纠错和必要弃权可以在同一个解码规则中协同发挥作用。

\section{引言}\label{sec:intro}
一张生牛肉片图像中，烤面包占据了显眼的位置。面对这张图，模型在十次独立试答中均未正确识别主菜；加入弃权要求后，正常分支已经倾向于回答 \mbox{\texttt{UNKNOWN}}。然而，将相同要求同时加入视觉对比的参考分支，最终回复却变成 Toast。一个用于减少视觉幻觉的操作，反而覆盖了模型原有的谨慎判断。图\ref{fig:teaser}把这一问题与答案纠正放在一起：三条件解码既恢复了生牛肉片案例中的必要弃权，也将另一张图像的 apple strudel 纠正为 APPLE PIE。可靠的视觉问答需要同时保留这两类收益。

视觉对比解码用正常输入与信息削弱输入之间的分布差异修正生成。VCD 通过扰动图像形成参考，M3ID 通过去除图像估计语言先验\citep{leng2024vcd,favero2024m3id}。这些差分使答案更加依赖视觉信息。将弃权指令直接加到各分支，看起来能够进一步处理证据不足的输入，但相同的指令会随视觉条件产生不同的响应。当参考图像的信息较弱时，弃权要求反而更容易得到满足；扣除参考分数也就扣除了这部分弃权优势。与此同时，正常分布在对比之前截断候选，落在集合之外的弃权词元已经失去成为输出的机会。

这两个操作把问题从“提示怎样写”推进到“指令怎样参与对比”。我们在相同回复前缀上分别控制视觉信息与弃权指令，得到四路分布。正常分布与参考分布的指令增量之差，直接显示对比解码额外抑制了哪些候选。再将这一局部变化与原有弃权的去向联系起来，就能区分答案纠正与错误替代。带引导的视觉对比在此承担机制对照的角色，原生方法仍是整体效果的比较基线。

据此，我们让带引导的正常分布提供任务要求和候选集合，让两条无弃权引导的分布专门计算视觉差分。三条件解码恰好保留一次正常视觉条件中的相对指令增量，同时保留由视觉证据决定的校正方向。它回答了对比解码中的一个具体设计问题：任务要求应当进入回答基础，视觉差分应当比较相同任务文本下的信息变化。

本文的贡献围绕这一认识展开。首先，我们发现并解释了视觉对比中的弃权流失：候选截断决定弃权能否进入解码，参考指令交互决定其进入后是否被抵消。其次，我们给出指令保持的三条件构造，建立分数分解、候选优势比与实际回复之间的联系。最后，两种参考方式、九个模型和两类视觉任务显示，这一构造能够同时改善答案纠正与必要弃权的联合决策表现。

\begin{figure*}[t]
\centering\includegraphics[width=\textwidth]{figures/fig01_authentic_case_teaser.pdf}
\caption{\textbf{保留必要弃权，同时纠正具体答案。} 左：LLaVA-Mistral 在生牛肉片输入上由双侧引导 VCD 的 Toast 转为 IP-VCD 的 \mbox{\texttt{UNKNOWN}}；底部标尺显示实际首词元的弃权边际跨过零点。右：Qwen3.5 的回复由 apple strudel 转为 APPLE PIE。图像与回复来自实际输入和生成记录。}
\label{fig:teaser}
\end{figure*}

\FloatBarrier
\section{相关工作}
\subsection{对比分布怎样控制生成}
VCD、M3ID 与 SID 分别利用扰动图像、去图像条件和内部视觉词元筛选构造参考\citep{leng2024vcd,favero2024m3id,huo2025sid}。它们通过相减提高正常视觉证据的相对作用。参考输入同时包含语言先验与任务文本，因此分布差分也会受到文本要求的影响。本文关注的正是这层耦合：参考分支对弃权指令的响应会随视觉条件变化，并直接进入最终分数。

DoLa 和 DeCo 利用层间预测或中间层信息校正输出，为解码期改进提供了另一类比较对象\citep{chuang2024dola,wang2025deco}。DExperts 将基础分布与专家、反专家分布组合，说明解码分数可以由不同职责的分布共同构成\citep{liu2021dexperts}。三条件解码沿用分布组合的形式，把正常任务引导与无引导视觉差分明确分开，从而确定指令作用在视觉校正中的系数。

\subsection{回答质量与必要弃权}
选择性预测把答案质量与作答范围共同纳入评价。Reliable VQA 用风险与覆盖率描述可靠作答，ReCoVERR 通过追加视觉问题收集证据，减少不必要的弃权\citep{whitehead2022reliablevqa,srinivasan2024recoverr}。CDA 根据知识相关性动态组合回答与弃权分支\citep{kim2025cda}。本文采用其视觉迁移实现，与原生解码和三条件方法比较联合决策效用。

不同的弃权参考对应不同的任务困难。Food-101 的独立试答和类别排序刻画指定模型能否识别主菜；VizWiz 的官方标签刻画用户问题能否从图像中回答。两项任务共同要求：答案有依据时完成作答，依据不足时保留弃权。将两种成功放在全部输入的同一分母下，能够直接判断增加的弃权是否带来净收益。

\subsection{从输出变化定位解码作用}
解码干预可能同时纠正部分错误并引入另一些错误，聚合准确率掩盖了这些变化的去向\citep{yin2025mirage,lee2024delve}。TRAPSBench 进一步显示，可回答性的内部表征与最终表达之间存在差距\citep{pramono2026trapsbench}。这些观察促使我们沿生成过程区分候选资格、相对分数和完整行为。共享前缀测量固定了当前生成位置，配对回复则揭示局部调整最终带来了正确作答、必要弃权还是新的错误。

\begin{figure*}[t]
\centering\includegraphics[width=\textwidth]{figures/fig02_mechanism_four_panel.pdf}
\caption{\textbf{弃权流失的行为后果与分数机制。} (a) LLaVA-Mistral 的287个原有合理弃权在两种解码下的去向。(b) Qwen2.5-VL 的主指令与参考指令交叉矩阵，颜色为语义弃权率，每格2,424题。(c) 生牛肉片案例中 UN 与 To 的四路对数优势及指令增量。(d) 同一前缀上连续增加参考指令交互，弃权边际在 $\eta=9/13$ 处翻转；曲线为实际四路分数的解析重组。}
\label{fig:mechanism}
\end{figure*}

\newpage
\section{弃权怎样在视觉对比中流失}\label{sec:mechanism}
\subsection{原有合理弃权被改写为错误}\label{sec:flow}
我们用独立试答确定 Food 的模型相对困难集合：正确类别在全类别排序中未列第一，且十次独立作答均未命中。它与带引导 Direct 的弃权取交集，得到具有试答依据的原有合理弃权。固定这个集合后，回复变化有三种去向：继续弃权、纠正为正确答案、改写为错误答案。

带引导 VCD 大量产生了第三种变化。LLaVA-Mistral 的共同 \mbox{\texttt{UNKNOWN}} 条件把原有合理弃权中的81.88\%改写为错误答案；IP-VCD 则使错误替代减少30例，并恢复更多必要弃权（图\ref{fig:mechanism}a）。对比解码因而需要处理的不只是答案竞争，还包括怎样保留已经获得支持的谨慎决定。

参考指令的交叉控制显示，这种流失与参考侧文本直接相关。保持正常图像和主指令不变，只替换参考表达，Qwen2.5-VL 的弃权率便大幅变化（图\ref{fig:mechanism}b）。四个固定主指令行的跨度均超过50个百分点。正常分支提供了相同的回答要求，最终行为却受参考文本强烈牵动，说明弃权指令已经进入本应提取视觉证据的差分项。

\begin{figure}[t]
\centering\includegraphics[width=\columnwidth]{figures/fig03_three_condition_duty_schematic.pdf}
\caption{\textbf{任务引导与视觉差分分工。} 带引导正常条件 $g$ 提供基础分数和候选集合 $\mathcal S_g$；无引导条件 $c,r$ 提供视觉差分。底部同名输入接收上方的 $g$ 与 $\mathcal S_g$。三路共享模型和生成前缀，组合分数在 $\mathcal S_g$ 上归一化。}
\label{fig:architecture}
\end{figure}

\subsection{候选资格与参考响应依次起作用}
固定图像 $x$、问题 $q$ 和实际回复前缀 $y_{<t}$。记无弃权引导的正常与参考对数概率为 $c,r$，加入相同弃权指令后的分布为 $g,h$。四路保留相同任务问题与简洁作答要求。正常与参考条件的指令增量分别为 $\delta_c=g-c$、$\delta_r=h-r$。原生 VCD 使用
\begin{equation}\begin{aligned}
s_N&=(1+\alpha)c-\alpha r,\\
\mathcal S_c&=\{v:c(v)\ge\max_u c(u)+\log\beta\}.
\end{aligned}\label{eq:native}\end{equation}
候选集合先决定哪些词元能被输出。无论视觉差分给出多大的相对增益，集合外的词元都不会参与最终选择。在九模型的每类一图代表面板中，指定弃权首词元均未进入无引导支持；加入弃权指令后，部分词元获得了候选资格。这个变化解释了引导正常分布对弃权的第一层作用。

获得候选资格后，参考侧的响应继续改变其排序。双侧引导 VCD 可展开为
\begin{equation}\begin{aligned}
s_G&=(1+\alpha)g-\alpha h\\
&=s_N+\delta_c+\alpha(\delta_c-\delta_r).
\end{aligned}\label{eq:guided}\end{equation}
对弃权候选 $u$ 与答案候选 $v$，令 $\Delta_c=\delta_c(u)-\delta_c(v)$、$\Delta_r=\delta_r(u)-\delta_r(v)$。当参考侧对弃权的相对响应更强，即 $\Delta_r>\Delta_c$，最后一项就降低弃权相对答案的分数。共同进入引导支持的候选对中，近四分之三呈现这一方向。引导让弃权获得竞争资格，参考交互又削弱其竞争优势，两者发生在不同步骤。

\subsection{分数交互能够翻转实际回答}
生牛肉片案例把上述分解接到完整回复。正常指令和参考指令都强烈提高 UN 相对 To 的优势，但参考增量更大（图\ref{fig:mechanism}c）。这部分额外响应在相减时压低弃权边际，使双侧引导 VCD 选择 To，并继续生成 Toast。移除额外交互后，边际由 $-0.1875$ 转为 $+0.4219$，实际回复转为 \mbox{\texttt{UNKNOWN}}。

参考交互的作用还可以连续分解。以保留一次正常指令增量的分数为起点，定义
\[s_\eta=s_N+\delta_c+\eta\alpha(\delta_c-\delta_r).\]
该案例的候选边际满足 $m_\eta=0.421875-0.609375\eta$，在 $\eta=9/13$ 时跨过零点（图\ref{fig:mechanism}d）。同一组实际分布由此连接了指令增量、候选次序和生成首词元：参考响应的额外扣除足以把有依据的弃权推向具体错误答案。

\section{指令保持的三条件解码}
\subsection{让任务要求留在回答基础中}
两阶段机制直接给出构造原则：弃权指令应决定回答基础及候选资格，视觉对比应在相同的无引导任务文本下提取证据差异。IP-VCD 据此使用
\begin{equation}\begin{aligned}
s_{\mathrm{IP}}&=g+\alpha(c-r),\\
\mathcal S_g&=\{v:g(v)\ge\max_u g(u)+\log\beta\}.
\end{aligned}\label{eq:ip}\end{equation}
输出分布为
\begin{equation}
p_{\mathrm{IP}}(v)=\frac{\mathbf1\{v\in\mathcal S_g\}e^{g(v)+\alpha[c(v)-r(v)]}}{\sum_{u\in\mathcal S_g}e^{g(u)+\alpha[c(u)-r(u)]}}.
\label{eq:prob}\end{equation}
引导正常分布首先表达怎样回答，视觉差分随后调整各候选获得的图像支持。改变视觉强度 $\alpha$ 时，任务指令的系数始终保持为1。三路由同一个冻结模型计算，图\ref{fig:architecture}展示其职责与数据流。

这个分工同时保留了答案纠正的路径。若正常图像比参考图像更支持某个具体答案，$c-r$ 仍会提高它的相对分数；若正常条件在困难输入上形成弃权倾向，$g$ 则将该要求带入最终分布。两类作用通过各自的条件进入生成，在同一候选集合中共同决定回复。

\begin{table*}[t]\centering
\caption{\textbf{Food-101 与 VizWiz 的联合决策比较。} 五个代表模型；指标单位为百分比。Food 每行 $N=2,424$，VizWiz 每行 $N=512$。IP 的 Food 指标取同一观测最高 $J$ 配置，VizWiz 为单一工作点。粗体为各模型、各任务的最高 $J$。}
\label{tab:main}
\fontsize{7.4}{9.0}\selectfont\setlength{\tabcolsep}{4.0pt}
\begin{tabular}{llrrrr@{\hspace{11pt}}rrrr}\toprule
模型 & 方法 & \multicolumn{4}{c}{Food-101} & \multicolumn{4}{c}{VizWiz}\\
\cmidrule(lr){3-6}\cmidrule(lr){7-10}
 & & Acc & P & R & $J$ & Q & P & R & $J$\\\midrule
Qwen2.5-VL & Direct & 30.40 & -- & 0.00 & 30.40 & 26.27 & 92.86 & 23.49 & 33.89\\
 & VCD & 28.34 & -- & 0.00 & 28.34 & 24.59 & 91.30 & 12.65 & 28.69\\
 & DoLa & 24.50 & -- & 0.00 & 24.50 & 20.29 & 86.49 & 19.28 & 26.54\\
 & DeCo & 30.24 & -- & 0.00 & 30.24 & 23.07 & 87.10 & 16.27 & 28.34\\
 & SID & 26.40 & -- & 0.00 & 26.40 & n/a & n/a & n/a & n/a\\
 & CDA & 33.91 & 23.54 & 23.24 & 37.58 & 13.12 & 41.77 & 59.64 & 32.46\\
 & IP-VCD & 36.30 & 26.12 & 63.71 & 46.37 & 14.00 & 52.19 & 78.92 & \textbf{39.59}\\
 & IP-M3ID & 42.16 & 25.31 & 37.34 & \textbf{48.06} & 11.89 & 53.20 & 80.12 & 37.87\\
\addlinespace[3pt]
LLaVA-Mistral & Direct & 33.87 & 77.78 & 0.74 & 34.16 & 19.37 & 85.71 & 25.30 & 27.58\\
 & VCD & 38.24 & -- & 0.00 & 38.24 & 16.35 & 83.33 & 6.02 & 18.30\\
 & DoLa & 33.29 & 71.43 & 0.53 & 33.50 & 18.91 & 83.72 & 21.69 & 25.94\\
 & DeCo & 26.90 & 50.00 & 0.11 & 26.94 & 12.70 & 85.71 & 18.07 & 18.55\\
 & SID & 38.28 & -- & 0.00 & 38.28 & 16.09 & 87.50 & 16.87 & 21.56\\
 & CDA & 31.31 & 50.00 & 0.42 & 31.48 & 15.47 & 51.43 & 32.53 & 26.02\\
 & IP-VCD & 28.38 & 53.04 & 61.54 & \textbf{52.48} & 14.51 & 53.88 & 79.52 & \textbf{40.29}\\
 & IP-M3ID & 20.17 & 46.52 & 61.22 & 44.14 & 10.70 & 42.98 & 90.36 & 40.00\\
\addlinespace[3pt]
MiniCPM-V2.6 & Direct & 43.03 & 100.00 & 0.50 & 43.15 & 29.18 & 94.74 & 10.84 & 32.70\\
 & VCD & 43.61 & 100.00 & 0.17 & 43.65 & 25.18 & 87.50 & 4.22 & 26.54\\
 & DoLa & 31.11 & 100.00 & 0.17 & 31.15 & 22.36 & 72.73 & 9.64 & 25.49\\
 & DeCo & 41.96 & 100.00 & 0.50 & 42.08 & 29.36 & 95.83 & 13.86 & 33.85\\
 & SID & n/a & n/a & n/a & n/a & n/a & n/a & n/a & n/a\\
 & CDA & 19.60 & -- & 0.00 & 19.60 & 12.54 & 49.72 & 54.22 & 30.12\\
 & IP-VCD & 37.25 & 46.74 & 42.04 & \textbf{47.61} & 26.52 & 64.04 & 68.67 & 48.79\\
 & IP-M3ID & 31.27 & 38.24 & 56.11 & 45.09 & 26.35 & 56.50 & 75.90 & \textbf{50.96}\\
\addlinespace[3pt]
InternVL3.5 & Direct & 17.08 & -- & 0.00 & 17.08 & 24.39 & 94.74 & 10.84 & 27.91\\
 & VCD & 19.10 & 100.00 & 0.09 & 19.14 & 22.64 & 92.86 & 7.83 & 25.18\\
 & DoLa & 16.17 & -- & 0.00 & 16.17 & 22.19 & 100.00 & 4.82 & 23.75\\
 & DeCo & 17.95 & -- & 0.00 & 17.95 & 21.39 & 100.00 & 6.63 & 23.54\\
 & SID & 18.52 & -- & 0.00 & 18.52 & 23.59 & 92.31 & 7.23 & 25.94\\
 & CDA & 20.38 & 100.00 & 0.09 & 20.42 & 19.96 & 68.49 & 30.12 & 29.73\\
 & IP-VCD & 27.72 & 82.50 & 3.07 & 29.08 & 16.43 & 69.42 & 50.60 & 32.83\\
 & IP-M3ID & 27.48 & 85.45 & 4.37 & \textbf{29.41} & 17.13 & 69.85 & 57.23 & \textbf{35.68}\\
\addlinespace[3pt]
Qwen3-VL & Direct & 53.92 & 50.00 & 0.27 & 53.96 & 37.52 & 96.43 & 16.27 & 42.79\\
 & VCD & 56.97 & -- & 0.00 & 56.97 & 35.37 & 96.15 & 15.06 & 40.25\\
 & DoLa & 51.82 & 100.00 & 0.27 & 51.86 & 36.23 & 83.87 & 15.66 & 41.31\\
 & DeCo & 53.71 & -- & 0.00 & 53.71 & 37.17 & 92.59 & 15.06 & 42.05\\
 & SID & 48.14 & -- & 0.00 & 48.14 & n/a & n/a & n/a & n/a\\
 & CDA & 2.76 & 0.00 & 0.00 & 2.76 & 17.52 & 62.09 & 57.23 & 36.07\\
 & IP-VCD & 53.01 & 36.80 & 46.70 & \textbf{60.02} & 27.21 & 66.86 & 70.48 & 50.06\\
 & IP-M3ID & 49.83 & 37.78 & 41.21 & 56.02 & 27.64 & 59.45 & 77.71 & \textbf{52.83}\\
\bottomrule\end{tabular}
\par\vspace{3pt}\begin{minipage}{.98\textwidth}\fontsize{7.2}{9}\selectfont
Acc：全部输入正确率；Q：答案共识信用；P/R：弃权精确率/召回率。--：零弃权时精确率未定义；n/a：原 SID 不适用。MiniCPM 的视觉词元不足100，Qwen3.5 第二层为线性注意力；VizWiz 中 Qwen2.5-VL、Qwen3-VL 各有7题少于固定100视觉词元。CDA 为视觉迁移。
\end{minipage}\end{table*}


\subsection{精确保留一次正常指令增量}
由式\eqref{eq:ip}可得 $s_{\mathrm{IP}}=s_N+\delta_c$。在固定前缀和共同支持内，任意两个候选满足
\begin{equation}\begin{aligned}
&[s_{\mathrm{IP}}(u)-s_{\mathrm{IP}}(v)]-[s_N(u)-s_N(v)]\\
&\hspace{1.2em}=\delta_c(u)-\delta_c(v).
\end{aligned}\label{eq:preserve}\end{equation}
归一化常数在优势比中抵消，正常视觉条件对候选间的相对指令作用因此被完整保留。相对于双侧引导 VCD，三条件构造加入 $\alpha(\delta_r-\delta_c)$，正好抵消式\eqref{eq:guided}中的额外交互。候选支持的变化由 $\mathcal S_g$ 负责，支持内的指令保持由这一恒等式负责。

第三条分布 $c$ 是这种分离所需的参照。只移除参考引导会得到 $(1+\alpha)g-\alpha r=s_N+(1+\alpha)\delta_c$，正常指令仍随视觉强度一起放大。加入无引导正常条件后，视觉差分与指令增量的系数才分别确定。在候选无关系数的线性族 $ag+bc+dr$ 中，同时要求单位指令响应和 $g=c$ 时恢复原生 VCD 候选边际，得到唯一系数组 $(1,\alpha,-\alpha)$。

三条件分布还具有清楚的优化解释。固定非空支持 $S=\mathcal S_g$，令 $G$ 为引导分布在 $S$ 内的归一化，式\eqref{eq:prob}唯一最大化
\begin{equation}
\alpha\,\mathbb E_\pi[c-r]-D_{\mathrm{KL}}(\pi\Vert G).
\label{eq:kl}\end{equation}
视觉差分奖励获得更多正常图像支持的候选，KL 项使偏移以引导分布为参照。分布的基础与偏移方向分别来自任务要求和视觉证据。

\subsection{两种参考方式遵循相同分工}
IP-VCD 使用扰动图像参考，设置 $\alpha=1,\beta=0.1$。IP-M3ID 使用去图像参考，并保留 M3ID 的位置相关强度与置信度激活：
\begin{equation}s_{\mathrm{IPM},t}=g_t+w_t^c(c_t-r_t),\label{eq:ipm}\end{equation}
\begin{equation}\begin{aligned}
w_t^c&=\mathbf1\{\max_v e^{c_t(v)}<\kappa\}(e^{\lambda\tau_t^c}-1),\\
\tau_t^c&=(t-1)+k_0^c.
\end{aligned}\label{eq:weight}\end{equation}
其中 $\lambda=0.02,\kappa=0.3$，$k_0^c$ 为无引导任务文本的词元长度，IP-M3ID 在完整词表上归一化。视觉强度由无引导条件确定，指令通过 $g_t$ 加入。对相同权重的原生分数，式\eqref{eq:preserve}中的单位指令关系继续成立。

自回归生成时，每一步把已选词元同时追加到三路缓存。相同前缀使各路分布比较同一待生成位置，独立缓存保留各条件的输入差异。这个操作在终止词元或32词元上限处结束。三条件规则直接作用于下一词元分布，所需的新增计算来自无引导正常条件。

\begin{figure*}[t]
\centering\includegraphics[width=\textwidth]{figures/fig04_food_nine_model_J_heatmap.pdf}
\caption{\textbf{九模型的 Food 联合效用。} 每模型2,424题。IP 行使用四个表达中的观测最高 $J$，其余方法各用一个登记工作点；CDA 为视觉迁移。金色外框标出三条件方法，粗体标出每行最高值。SID 的架构适用性由表注说明。}
\label{fig:food}
\end{figure*}

\section{回答与必要弃权的联合收益}
\subsection{任务与联合效用}
Food-101 每个条件包含101类各24张图像，以主菜的规范类别名称判定正确性\citep{bossard2014food101}；弃权参考为第\ref{sec:flow}节的模型相对困难集合。VizWiz 使用固定512题，包含346个官方可回答输入与166个不可回答输入，回答质量取官方连续共识得分\citep{gurari2018vizwiz,vizwizofficial}。完整回复按语义区分作答与弃权。九个模型使用相同任务输入，涵盖 Qwen、LLaVA、MiniCPM、Gemma、InternVL、OneVision 和 Phi 系列。

联合效用把正确作答与有参考支持的弃权共同计为成功。令 $C,W$ 为正确、错误作答数，$TP,FP$ 为参考阳性、阴性弃权数，则
\begin{equation}J_{\mathrm{Food}}=\frac{C+TP}{N}=1-\frac{W+FP}{N}.\label{eq:foodj}\end{equation}
这一指标直接衡量正确处理了多少输入，增加无依据弃权会与错误作答一样受到惩罚。VizWiz 将二元正确性换成答案共识信用 $q_i\in[0,1]$。令 $a_i$ 为语义弃权、$u_i$ 为官方不可回答标签，则
\begin{equation}\begin{aligned}
Q&=\frac1N\sum_i(1-a_i)q_i,\\
J_{\mathrm{Viz}}&=Q+\frac1N\sum_i a_i u_i.
\end{aligned}\label{eq:vizj}\end{equation}
同时报告整体准确率或回答信用，以及弃权精确率 $P=TP/(TP+FP)$ 和召回率 $R=TP/N_+$，以解释联合收益来自哪种决策变化。

主比较包括原生 Direct、VCD、DoLa、DeCo、适用架构上的 SID，以及 CDA 视觉迁移和两个 IP 实例。Food 的 IP 行取四个预定表达中观测最高 $J$ 的单一配置，全部指标来自同一行；基线各用一个登记工作点。VizWiz 使用各方法的单一记录工作点。带引导 VCD 和 M3ID 用于相同指令下的机制消融。各条件均采用冻结模型、贪心生成和32词元预算。

\begin{figure}[t]
\centering\includegraphics[width=\columnwidth]{figures/fig05_vizwiz_paired_J_forest.pdf}
\caption{\textbf{VizWiz 上相对原生 VCD 的联合收益。} 每个点为相同512题上的配对 $\Delta J$，误差线为按输入重采样得到的95\%百分位 bootstrap 区间，2,000次重采样。圆点和菱形分别表示 IP-VCD 与 IP-M3ID；虚线为原生 VCD 基准。}
\label{fig:viz}
\end{figure}

\subsection{视觉纠错与谨慎作答共同提高收益}
三条件构造在五个代表模型上取得最高联合效用（表\ref{tab:main}），但收益并不依赖单一的行为模式。Qwen2.5-VL 与 InternVL 的提升同时表现为更多正确答案，说明视觉差分继续承担细粒度识别；LLaVA 和 MiniCPM 的收益更多来自困难输入上的必要弃权。把这两条路径纳入同一个指标，才能判断解码是否真正改善了最终决定。

Food 的关键变化是：难以识别的输入不再只能在具体类别之间竞争，具有试答依据的弃权也能成为输出。以 Qwen2.5-VL 为例，IP-VCD 相对原生 VCD 的联合效用提高18.03个百分点。增加的必要弃权覆盖了原生对比继续作答的困难输入，答案纠正则进一步提高成功处理的比例。精确率与召回率保留了两类弃权的区别，使净收益直接对应正确答案和有依据的拒答。

九模型全景显示，这种收益具有广泛的架构覆盖（图\ref{fig:food}）。IP-VCD 在八个模型上超过原生 VCD，IP-M3ID 在六个模型上提高联合效用；两个实例组成的方法组在八个模型上取得所比较方法的最高值。不同参考方式能够利用不同模型的视觉与语言响应，同一三条件原则因而表现为互补的决策改进。

CDA 对照进一步显示，主动引入弃权分支与保持任务指令具有不同作用。其视觉迁移通过熵相关权重同时调整知识与弃权分支；实测轨迹包含负弃权系数，部分输入上的分支权重幅度明显放大。动态线性组合由此可以削弱某些原有的弃权响应。三条件构造把指令系数固定为1，使视觉证据的强度变化与任务要求的保留分别受到控制。

\subsection{真实不可回答问题上的改善更一致}
VizWiz 把模型相对困难推进到图像本身能否回答用户问题。模糊视野、缺失内容和问题所需信息不可见，都要求模型在生成时识别证据不足。两个 IP 实例在九模型上均提高相对原生 VCD 的联合效用（图\ref{fig:viz}）。在包含 DoLa、DeCo、适用架构的 SID 与 CDA 的完整比较中，两个实例也分别超过全部登记基线；IP-VCD 与 IP-M3ID 分别取得四个和五个模型的最高联合效用（附录图\ref{fig:viz-all}）。这说明正常视觉条件中的弃权响应能够转化为实际不可回答问题上的成功决定。

改善同时包含答案质量和可回答性判断。MiniCPM 在维持并提高答案信用的同时，覆盖了更多不可回答输入；LLaVA 的主要收益来自减少缺乏依据的具体回答。不同模型因此获得了不同的回答与弃权组合，整体效用都向同一方向变化。表\ref{tab:main}中的精确率和召回率进一步刻画这些工作点，使必要弃权的增加与可回答输入上的弃权区别开来。

两项任务共同显示，弃权的价值取决于它是否发生在需要弃权的输入上。Food 的参考由同一模型的独立试答确定，VizWiz 的参考来自官方可回答性标签；三条件解码在两种定义下都使已有的视觉判断更容易保留到最终回复。联合效用将这种收益与答案质量一同计入，而完整的决策指标解释它如何发生。

\begin{table*}[t]\centering
\caption{\textbf{相同指令下的结构作用与推理开销。} 左：九模型、四表达的36个配对，$\Delta J$ 单位为百分点。右：同硬件、同表达的匹配运行，每方法2,424题；VCD 为同措辞引导条件，时间包含输入处理、会话构造与生成。}
\label{tab:ablation-cost}
\begin{minipage}[t]{.51\textwidth}\centering\fontsize{8}{10}\selectfont\setlength{\tabcolsep}{3pt}
\begin{tabular}{llrr}\toprule
实例 & 对照 & 平均 $\Delta J$ & 范围\\\midrule
IP-VCD & 双侧引导 & +0.48 & [$-$2.72, +2.52]\\
IP-VCD & 参考去引导 & +3.19 & [$-$8.87, +25.78]\\
IP-M3ID & 双侧引导 & $-$0.06 & [$-$4.74, +3.63]\\
IP-M3ID & 参考去引导 & $-$1.33 & [$-$7.55, +4.87]\\\bottomrule
\end{tabular}\end{minipage}\hfill
\begin{minipage}[t]{.47\textwidth}\centering\fontsize{8}{10}\selectfont\setlength{\tabcolsep}{3pt}
\begin{tabular}{llrrr}\toprule
模型 & 表达 & VCD & IP-VCD & CDA\\
 & & \multicolumn{3}{c}{秒/题（相对 VCD）}\\\midrule
Qwen2.5-VL & UNKNOWN & 0.706 & 1.051 & 1.833\\
 & & (1.00$\times$) & (1.49$\times$) & (2.60$\times$)\\
MiniCPM-V2.6 & UNKNOWN & 0.848 & 1.322 & 2.052\\
 & & (1.00$\times$) & (1.56$\times$) & (2.42$\times$)\\
\bottomrule\end{tabular}\end{minipage}
\end{table*}


\subsection{候选支持决定指令保持发挥作用的位置}
Gemma-3 揭示了两阶段机制中的前置条件。它在 Food 的代表面板上没有将指定弃权词元纳入引导支持，因此该任务中的 IP-VCD 联合效用略低于原生 VCD；IP-M3ID 也略低于 Direct。相同模型在 VizWiz 中表现出更多有效弃权，并获得联合收益。任务改变了引导正常分布中的候选基础，分数修正能够发挥的空间也随之改变。

这一差异把模型间结果接回候选操作。指令保持保证共同支持内的相对增量，实际可选词元仍由引导正常分布决定。引导能够使弃权进入支持时，视觉差分与任务要求共同参与竞争；引导支持集中在具体答案时，方法的作用主要体现在答案间的校正。这一条件同时解释了不同模型和不同任务的响应差异。

\section{结构消融与计算开销}
\subsection{同一指令下恢复必要弃权}
相同主指令与候选支持使局部分数交互得到直接比较。LLaVA 的 \mbox{\texttt{UNKNOWN}} 条件中，IP-VCD 与双侧引导 VCD 的总正确数相同，但原有合理弃权的错误替代减少了30例（图\ref{fig:mechanism}a）。改变来自参考指令交互的移除，说明三条件构造能够在保持答案总正确数的同时恢复必要弃权。

参考去引导消融进一步检验第三个条件 $c$ 的作用。双条件变体保留 $(1+\alpha)\delta_c$，三条件将正常指令增量固定为 $\delta_c$，两者因此具有不同的指令强度。九模型的同表达比较中，IP-VCD 相对参考去引导变体平均提高联合效用（表\ref{tab:ablation-cost}）。独立控制指令系数有助于减少被视觉强度共同放大的不必要弃权。

IP-M3ID 的同指令变化还包含置信度门控与位置强度的差异，配对结果同时呈现提升与下降。它在主比较中的整体收益来自引导基础与无图像参考的共同作用；局部消融则揭示强度调度怎样影响这条路径。两种实例由此共享三条件原则，同时保留各自参考构造的行为特点。

\subsection{三路状态带来明确的推理代价}
IP-VCD 维护三路条件状态，原生 VCD 维护两路，CDA 视觉迁移使用五路。实际耗时还取决于回复长度、视觉编码与缓存。我们比较14个同模型、同表达的完整匹配组，每组各方法均有2,424题，组内硬件与软件身份一致。IP-VCD 的每题时延为匹配引导 VCD 的1.34--1.72倍，CDA 为2.42--3.57倍。

这些实测值将计算结构与运行开销联系起来。三条件计算以新增一路正常分布，换取任务引导和视觉差分的独立控制；相对于五路 CDA，分支数量与实测延迟均更低。表\ref{tab:ablation-cost}给出匹配工作点的具体时延。计时涵盖图像读取、会话构造、生成及清理，模型加载在计时区间之外。

\newpage
\section{结论}
视觉对比中的弃权流失源于两个相继发生的操作：候选集合先决定弃权是否可选，参考指令交互再改变它与具体答案的相对优势。三条件解码将任务引导保留在正常回答基础中，由无引导条件单独提取视觉差分，从而同时承接视觉纠错与必要弃权。四路分布与真实回复把这一构造的作用具体化，两类任务的联合收益则显示其决策价值。可靠的多模态对比生成因此需要同时考虑比较何种证据，以及保留何种回答要求。

\section*{适用范围与伦理}
Food 的参考刻画固定试答协议下的模型作答困难，VizWiz 的参考刻画输入的官方可回答性。Food 四表达观测最优值描述有限候选集合的经验表现；局部分数恒等式对应固定前缀和共同支持。实际应用可按错误作答与不必要弃权的代价设定效用权重。实验使用公开数据与冻结模型；VizWiz 中真实用户的图像和问题以必要的汇总统计呈现。

\clearpage\nobalance
\small\setlength{\bibsep}{2pt}\bibliographystyle{unsrtnat}\bibliography{refs/references}
\clearpage\onecolumn\renewcommand{\baselinestretch}{1.08}\normalsize
\appendix\section{任务与运行设置}\label{ux4efbux52a1ux4e0eux8fd0ux884cux8bbeux7f6e}

\subsection{数据与参考集合}\label{ux6570ux636eux4e0eux53c2ux8003ux96c6ux5408}

Food-101 每条件包含101类各24题，共2,424题。类别评分先从实际回答确定明确主菜，再与101个规范名称匹配，统一大小写及首尾排版格式，接受下划线转空格形式；同义词、词形变化与拼写修复不加入主评分。配菜与解释性提及不作为主菜，多项同级或相互竞争的具体主菜按单标签任务记零。明确表达无法识别的完整回复计为语义弃权；仍给出具体答案的犹疑措辞按作答处理。

模型级参考集合同时满足正确类别的闭集排名大于1、十次独立最佳估计试答均失败。闭集名称分数采用平均词元对数概率。参考集合在该模型的方法比较中固定，方法实际答对的样本仍计入正确作答。各模型分母如下。

\begin{longtable}[]{@{}ll@{}}
\toprule\noalign{}
模型 & 参考阳性数 \\
\midrule\noalign{}
\endhead
\bottomrule\noalign{}
\endlastfoot
Gemma-3 & 422 \\
InternVL3.5 & 1075 \\
LLaVA-Mistral & 949 \\
MiniCPM-V2.6 & 597 \\
OneVision & 814 \\
Phi-3.5-Vision & 856 \\
Qwen2.5-VL & 383 \\
Qwen3.5 & 486 \\
Qwen3-VL & 364 \\
\end{longtable}

VizWiz 的固定512题包括346题官方可回答与166题不可回答。回答质量来自已标注的 answer\_text 经官方共识评分后的结果；语义弃权的主质量为零。原始整段回复的 official\_raw\_score 作为另一列保留。主质量与原始整段回复评分使用不同文本表征，因此正文 Q 及 J 均明确使用前者。

\subsection{提示与配置选择}\label{ux63d0ux793aux4e0eux914dux7f6eux9009ux62e9}

无弃权引导条件保留任务问题及``Give a concise answer.''；引导条件加入``If you cannot determine the answer reliably, respond with \ldots''。登记表达为 UNKNOWN、UNCLEAR、UNSURE 与 I cannot identify it。每步所有分支接收相同实际生成前缀，采用独立输入状态和缓存。

Food 每个 IP 实例、每个模型从四个登记表达中取观测 J 最高的整行，其他指标随该行一起报告；若 J 并列，按登记表顺序取首行。基线各有一个主登记工作点，SID 覆盖七模型。完整表共70行，选中表达列保留实际输入身份。VizWiz 主比较共68个工作点：九模型的 Direct、VCD、DoLa、DeCo、CDA 与两个 IP 实例，以及五个适用模型的原 SID；每个工作点512题。

SID 在 Gemma-3 上采用作者贪心算子：第二层注意力选择100个视觉词元、$\alpha=0.5$、完整词表对比。其余六模型的全量结果使用原登记的$\alpha=1$与正常分布支持截断；同题101张的作者算子配对结果见附录\ref{sec:supp-author-operators}。MiniCPM-V2.6有857/2,424张图像只含64个视觉词元，低于原法固定选择的100个；Qwen3.5第二层为线性注意力，不产生SID所需的softmax注意力权重。

VizWiz 的 SID 使用原作者算子，覆盖 Gemma、InternVL、LLaVA-Mistral、OneVision 与 Phi；Qwen2.5-VL 和 Qwen3-VL 各有7/512题的视觉词元少于固定100个，因而整面板记为不适用。DoLa、DeCo 在九模型上使用本轮核对后的原生无引导实现。

\subsection{模型与解码参数}\label{ux6a21ux578bux4e0eux89e3ux7801ux53c2ux6570}

\begin{longtable}[]{@{}ll@{}}
\toprule\noalign{}
模型 & 冻结检查点 \\
\midrule\noalign{}
\endhead
\bottomrule\noalign{}
\endlastfoot
Qwen2.5-VL & Qwen/Qwen2.5-VL-7B-Instruct \\
Qwen3.5 & Qwen/Qwen3.5-4B \\
LLaVA-Mistral & llava-hf/llava-v1.6-mistral-7b-hf \\
MiniCPM-V2.6 & openbmb/MiniCPM-V-2\_6 \\
Gemma-3 & google/gemma-3-4b-it \\
InternVL3.5 & OpenGVLab/InternVL3\_5-8B \\
OneVision & llava-hf/llava-onevision-qwen2-7b-ov-hf \\
Phi-3.5-Vision & microsoft/Phi-3.5-vision-instruct \\
Qwen3-VL & Qwen/Qwen3-VL-8B-Instruct \\
\end{longtable}

全部生成使用温度0、top-p为1、最大32词元；模型权重、处理器和精度沿登记设置。IP-VCD 使用扰动图像参考，α=1、β=0.1。IP-M3ID 使用去图像参考，λ=0.02、κ=0.3，位置索引为生成步的零起始位置加该输入无引导任务文本的登记词元长度。该偏移因问题文本而变化。DeCo 的登记默认参数为α=0.6、top-k=20、top-p=0.9；具体层、视觉输入及 SID 设置使用模型原生适配的登记配置。

\begin{longtable}[]{@{}lrl@{}}
\toprule\noalign{}
方法 & 条件数 & 条件职责 \\
\midrule\noalign{}
\endhead
\bottomrule\noalign{}
\endlastfoot
Direct & 1 & 正常回答基础 \\
VCD & 2 & 正常视觉与扰动参考 \\
IP-VCD & 3 & 引导基础、无引导正常视觉、无引导扰动参考 \\
IP-M3ID & 3 & 引导基础、无引导正常视觉、无引导文本参考 \\
CDA 视觉迁移 & 5 & 知识条件、弃权条件与两路空输入校准 \\
\end{longtable}

CDA 视觉迁移按原论文动态知识相关性规则构造权重，不使用 momentum。权重轨迹允许弃权分支系数为负，故其分数是动态 logit 线性组合。M3ID 的数学定义和支持规则在独立推导中给出。

\section{完整九模型结果}\label{ux5b8cux6574ux4e5dux6a21ux578bux7ed3ux679c}

\subsection{Food 全部工作点}\label{food-ux5168ux90e8ux5de5ux4f5cux70b9}

每行均来自同一个实际登记条件。Acc、P、R、J 单位为百分比；TP 与 FP 为次数；所有行 N=2,424。表达``NONE''表示无弃权引导，CDA 展示其登记输入。空精确率表示该条件没有语义弃权。

\begingroup\fontsize{8.2}{10.5}\selectfont\setlength{\tabcolsep}{3.7pt}
\begin{longtable}{llrrrrrrl}
\toprule
模型 & 方法 & Acc & P & R & J & TP & FP & 表达\\
\midrule\endfirsthead
\toprule
模型 & 方法 & Acc & P & R & J & TP & FP & 表达\\
\midrule\endhead
\bottomrule\endfoot
Qwen2.5-VL & Direct & 30.40 & -- & 0.00 & 30.40 & 0 & 0 & NONE\\
Qwen2.5-VL & VCD & 28.34 & -- & 0.00 & 28.34 & 0 & 0 & NONE\\
Qwen2.5-VL & DoLa & 24.50 & -- & 0.00 & 24.50 & 0 & 0 & NONE\\
Qwen2.5-VL & DeCo & 30.24 & -- & 0.00 & 30.24 & 0 & 0 & NONE\\
Qwen2.5-VL & SID & 26.40 & -- & 0.00 & 26.40 & 0 & 0 & NONE\\
Qwen2.5-VL & CDA & 33.91 & 23.54 & 23.24 & 37.58 & 89 & 289 & UNKNOWN\\
Qwen2.5-VL & IP-VCD & 36.30 & 26.12 & 63.71 & 46.37 & 244 & 690 & I cannot identify it\\
Qwen2.5-VL & IP-M3ID & 42.16 & 25.31 & 37.34 & 48.06 & 143 & 422 & I cannot identify it\\
LLaVA-Mistral & Direct & 33.87 & 77.78 & 0.74 & 34.16 & 7 & 2 & NONE\\
LLaVA-Mistral & VCD & 38.24 & -- & 0.00 & 38.24 & 0 & 0 & NONE\\
LLaVA-Mistral & DoLa & 33.29 & 71.43 & 0.53 & 33.50 & 5 & 2 & NONE\\
LLaVA-Mistral & DeCo & 26.90 & 50.00 & 0.11 & 26.94 & 1 & 1 & NONE\\
LLaVA-Mistral & SID & 38.28 & -- & 0.00 & 38.28 & 0 & 0 & NONE\\
LLaVA-Mistral & CDA & 31.31 & 50.00 & 0.42 & 31.48 & 4 & 4 & UNKNOWN\\
LLaVA-Mistral & IP-VCD & 28.38 & 53.04 & 61.54 & 52.48 & 584 & 517 & I cannot identify it\\
LLaVA-Mistral & IP-M3ID & 20.17 & 46.52 & 61.22 & 44.14 & 581 & 668 & UNSURE\\
MiniCPM-V2.6 & Direct & 43.03 & 100.00 & 0.50 & 43.15 & 3 & 0 & NONE\\
MiniCPM-V2.6 & VCD & 43.61 & 100.00 & 0.17 & 43.65 & 1 & 0 & NONE\\
MiniCPM-V2.6 & DoLa & 31.11 & 100.00 & 0.17 & 31.15 & 1 & 0 & NONE\\
MiniCPM-V2.6 & DeCo & 41.96 & 100.00 & 0.50 & 42.08 & 3 & 0 & NONE\\
MiniCPM-V2.6 & CDA & 19.60 & -- & 0.00 & 19.60 & 0 & 0 & UNKNOWN\\
MiniCPM-V2.6 & IP-VCD & 37.25 & 46.74 & 42.04 & 47.61 & 251 & 286 & UNCLEAR\\
MiniCPM-V2.6 & IP-M3ID & 31.27 & 38.24 & 56.11 & 45.09 & 335 & 541 & UNCLEAR\\
InternVL3.5 & Direct & 17.08 & -- & 0.00 & 17.08 & 0 & 0 & NONE\\
InternVL3.5 & VCD & 19.10 & 100.00 & 0.09 & 19.14 & 1 & 0 & NONE\\
InternVL3.5 & DoLa & 16.17 & -- & 0.00 & 16.17 & 0 & 0 & NONE\\
InternVL3.5 & DeCo & 17.95 & -- & 0.00 & 17.95 & 0 & 0 & NONE\\
InternVL3.5 & SID & 18.52 & -- & 0.00 & 18.52 & 0 & 0 & NONE\\
InternVL3.5 & CDA & 20.38 & 100.00 & 0.09 & 20.42 & 1 & 0 & UNKNOWN\\
InternVL3.5 & IP-VCD & 27.72 & 82.50 & 3.07 & 29.08 & 33 & 7 & I cannot identify it\\
InternVL3.5 & IP-M3ID & 27.48 & 85.45 & 4.37 & 29.41 & 47 & 8 & I cannot identify it\\
Qwen3-VL & Direct & 53.92 & 50.00 & 0.27 & 53.96 & 1 & 1 & NONE\\
Qwen3-VL & VCD & 56.97 & -- & 0.00 & 56.97 & 0 & 0 & NONE\\
Qwen3-VL & DoLa & 51.82 & 100.00 & 0.27 & 51.86 & 1 & 0 & NONE\\
Qwen3-VL & DeCo & 53.71 & -- & 0.00 & 53.71 & 0 & 0 & NONE\\
Qwen3-VL & SID & 48.14 & -- & 0.00 & 48.14 & 0 & 0 & NONE\\
Qwen3-VL & CDA & 2.76 & 0.00 & 0.00 & 2.76 & 0 & 1 & UNKNOWN\\
Qwen3-VL & IP-VCD & 53.01 & 36.80 & 46.70 & 60.02 & 170 & 292 & I cannot identify it\\
Qwen3-VL & IP-M3ID & 49.83 & 37.78 & 41.21 & 56.02 & 150 & 247 & UNCLEAR\\
Qwen3.5 & Direct & 51.32 & -- & 0.00 & 51.32 & 0 & 0 & NONE\\
Qwen3.5 & VCD & 51.90 & -- & 0.00 & 51.90 & 0 & 0 & NONE\\
Qwen3.5 & DoLa & 51.69 & -- & 0.00 & 51.69 & 0 & 0 & NONE\\
Qwen3.5 & DeCo & 51.16 & -- & 0.00 & 51.16 & 0 & 0 & NONE\\
Qwen3.5 & CDA & 26.28 & 0.65 & 0.21 & 26.32 & 1 & 152 & UNKNOWN\\
Qwen3.5 & IP-VCD & 51.90 & 46.81 & 4.53 & 52.81 & 22 & 25 & I cannot identify it\\
Qwen3.5 & IP-M3ID & 53.63 & 71.43 & 1.03 & 53.84 & 5 & 2 & I cannot identify it\\
Gemma-3 & Direct & 54.79 & -- & 0.00 & 54.79 & 0 & 0 & NONE\\
Gemma-3 & VCD & 53.71 & 100.00 & 0.24 & 53.75 & 1 & 0 & NONE\\
Gemma-3 & DoLa & 54.29 & -- & 0.00 & 54.29 & 0 & 0 & NONE\\
Gemma-3 & DeCo & 54.58 & -- & 0.00 & 54.58 & 0 & 0 & NONE\\
Gemma-3 & SID & 51.94 & -- & 0.00 & 51.94 & 0 & 0 & NONE\\
Gemma-3 & CDA & 30.61 & 32.00 & 1.90 & 30.94 & 8 & 17 & UNKNOWN\\
Gemma-3 & IP-VCD & 53.05 & 100.00 & 0.24 & 53.09 & 1 & 0 & UNCLEAR\\
Gemma-3 & IP-M3ID & 54.62 & -- & 0.00 & 54.62 & 0 & 0 & UNKNOWN\\
OneVision & Direct & 30.57 & -- & 0.00 & 30.57 & 0 & 0 & NONE\\
OneVision & VCD & 35.02 & -- & 0.00 & 35.02 & 0 & 0 & NONE\\
OneVision & DoLa & 18.94 & -- & 0.00 & 18.94 & 0 & 0 & NONE\\
OneVision & DeCo & 31.27 & -- & 0.00 & 31.27 & 0 & 0 & NONE\\
OneVision & SID & 33.75 & -- & 0.00 & 33.75 & 0 & 0 & NONE\\
OneVision & CDA & 25.87 & 100.00 & 0.12 & 25.91 & 1 & 0 & UNKNOWN\\
OneVision & IP-VCD & 36.39 & 75.00 & 1.11 & 36.76 & 9 & 3 & I cannot identify it\\
OneVision & IP-M3ID & 31.15 & 82.35 & 1.72 & 31.72 & 14 & 3 & I cannot identify it\\
Phi-3.5-Vision & Direct & 35.73 & -- & 0.00 & 35.73 & 0 & 0 & NONE\\
Phi-3.5-Vision & VCD & 37.79 & -- & 0.00 & 37.79 & 0 & 0 & NONE\\
Phi-3.5-Vision & DoLa & 30.61 & -- & 0.00 & 30.61 & 0 & 0 & NONE\\
Phi-3.5-Vision & DeCo & 29.08 & -- & 0.00 & 29.08 & 0 & 0 & NONE\\
Phi-3.5-Vision & SID & 35.89 & -- & 0.00 & 35.89 & 0 & 0 & NONE\\
Phi-3.5-Vision & CDA & 27.48 & 100.00 & 0.70 & 27.72 & 6 & 0 & UNKNOWN\\
Phi-3.5-Vision & IP-VCD & 36.88 & 69.23 & 5.26 & 38.74 & 45 & 20 & I cannot identify it\\
Phi-3.5-Vision & IP-M3ID & 33.50 & 48.31 & 5.02 & 35.27 & 43 & 46 & I cannot identify it\\
\end{longtable}\endgroup

\subsection{VizWiz 全部记录工作点}
\begin{figure}[htbp]\centering
\includegraphics[width=.90\textwidth]{figures/figA01_vizwiz_all_baselines_J.pdf}
\caption{\textbf{VizWiz 的完整九模型基线比较。} 固定512题，共68个工作点。金色框标出三条件方法，粗体为每模型最高联合效用；n/a为原SID的架构或视觉词元数量限制。}
\label{fig:viz-all}\end{figure}
\label{vizwiz-ux5168ux90e8ux8bb0ux5f55ux5de5ux4f5cux70b9}

Q 为行为调整后的答案文本共识信用，P 与 R 为弃权精确率、召回率；百分比列均保留两位小数。全部行 N=512，参考阳性为166。

\begingroup\fontsize{8.2}{10.5}\selectfont\setlength{\tabcolsep}{3.7pt}
\begin{longtable}{llrrrrrrl}
\toprule
模型 & 方法 & Q & P & R & J & TP & FP & 表达\\
\midrule\endfirsthead
\toprule
模型 & 方法 & Q & P & R & J & TP & FP & 表达\\
\midrule\endhead
\bottomrule\endfoot
Qwen2.5-VL & Direct & 26.27 & 92.86 & 23.49 & 33.89 & 39 & 3 & UNKNOWN\\
Qwen2.5-VL & VCD & 24.59 & 91.30 & 12.65 & 28.69 & 21 & 2 & NONE\\
Qwen2.5-VL & DoLa & 20.29 & 86.49 & 19.28 & 26.54 & 32 & 5 & NONE\\
Qwen2.5-VL & DeCo & 23.07 & 87.10 & 16.27 & 28.34 & 27 & 4 & NONE\\
Qwen2.5-VL & CDA & 13.12 & 41.77 & 59.64 & 32.46 & 99 & 138 & UNKNOWN\\
Qwen2.5-VL & IP-VCD & 14.00 & 52.19 & 78.92 & 39.59 & 131 & 120 & UNSURE\\
Qwen2.5-VL & IP-M3ID & 11.89 & 53.20 & 80.12 & 37.87 & 133 & 117 & I cannot identify it\\
LLaVA-Mistral & Direct & 19.37 & 85.71 & 25.30 & 27.58 & 42 & 7 & UNKNOWN\\
LLaVA-Mistral & VCD & 16.35 & 83.33 & 6.02 & 18.30 & 10 & 2 & NONE\\
LLaVA-Mistral & DoLa & 18.91 & 83.72 & 21.69 & 25.94 & 36 & 7 & NONE\\
LLaVA-Mistral & DeCo & 12.70 & 85.71 & 18.07 & 18.55 & 30 & 5 & NONE\\
LLaVA-Mistral & SID & 16.09 & 87.50 & 16.87 & 21.56 & 28 & 4 & NONE\\
LLaVA-Mistral & CDA & 15.47 & 51.43 & 32.53 & 26.02 & 54 & 51 & UNSURE\\
LLaVA-Mistral & IP-VCD & 14.51 & 53.88 & 79.52 & 40.29 & 132 & 113 & I cannot identify it\\
LLaVA-Mistral & IP-M3ID & 10.70 & 42.98 & 90.36 & 40.00 & 150 & 199 & I cannot identify it\\
MiniCPM-V2.6 & Direct & 29.18 & 94.74 & 10.84 & 32.70 & 18 & 1 & UNKNOWN\\
MiniCPM-V2.6 & VCD & 25.18 & 87.50 & 4.22 & 26.54 & 7 & 1 & NONE\\
MiniCPM-V2.6 & DoLa & 22.36 & 72.73 & 9.64 & 25.49 & 16 & 6 & NONE\\
MiniCPM-V2.6 & DeCo & 29.36 & 95.83 & 13.86 & 33.85 & 23 & 1 & NONE\\
MiniCPM-V2.6 & CDA & 12.54 & 49.72 & 54.22 & 30.12 & 90 & 91 & I cannot identify it\\
MiniCPM-V2.6 & IP-VCD & 26.52 & 64.04 & 68.67 & 48.79 & 114 & 64 & UNCLEAR\\
MiniCPM-V2.6 & IP-M3ID & 26.35 & 56.50 & 75.90 & 50.96 & 126 & 97 & UNCLEAR\\
InternVL3.5 & Direct & 24.39 & 94.74 & 10.84 & 27.91 & 18 & 1 & UNKNOWN\\
InternVL3.5 & VCD & 22.64 & 92.86 & 7.83 & 25.18 & 13 & 1 & NONE\\
InternVL3.5 & DoLa & 22.19 & 100.00 & 4.82 & 23.75 & 8 & 0 & NONE\\
InternVL3.5 & DeCo & 21.39 & 100.00 & 6.63 & 23.54 & 11 & 0 & NONE\\
InternVL3.5 & SID & 23.59 & 92.31 & 7.23 & 25.94 & 12 & 1 & NONE\\
InternVL3.5 & CDA & 19.96 & 68.49 & 30.12 & 29.73 & 50 & 23 & UNKNOWN\\
InternVL3.5 & IP-VCD & 16.43 & 69.42 & 50.60 & 32.83 & 84 & 37 & I cannot identify it\\
InternVL3.5 & IP-M3ID & 17.13 & 69.85 & 57.23 & 35.68 & 95 & 41 & I cannot identify it\\
Qwen3-VL & Direct & 37.52 & 96.43 & 16.27 & 42.79 & 27 & 1 & UNKNOWN\\
Qwen3-VL & VCD & 35.37 & 96.15 & 15.06 & 40.25 & 25 & 1 & NONE\\
Qwen3-VL & DoLa & 36.23 & 83.87 & 15.66 & 41.31 & 26 & 5 & NONE\\
Qwen3-VL & DeCo & 37.17 & 92.59 & 15.06 & 42.05 & 25 & 2 & NONE\\
Qwen3-VL & CDA & 17.52 & 62.09 & 57.23 & 36.07 & 95 & 58 & UNKNOWN\\
Qwen3-VL & IP-VCD & 27.21 & 66.86 & 70.48 & 50.06 & 117 & 58 & I cannot identify it\\
Qwen3-VL & IP-M3ID & 27.64 & 59.45 & 77.71 & 52.83 & 129 & 88 & UNCLEAR\\
Qwen3.5 & Direct & 37.15 & 100.00 & 0.60 & 37.34 & 1 & 0 & UNKNOWN\\
Qwen3.5 & VCD & 32.19 & 100.00 & 3.01 & 33.16 & 5 & 0 & NONE\\
Qwen3.5 & DoLa & 36.88 & 100.00 & 0.60 & 37.07 & 1 & 0 & NONE\\
Qwen3.5 & DeCo & 35.51 & 100.00 & 3.01 & 36.48 & 5 & 0 & NONE\\
Qwen3.5 & CDA & 22.21 & 49.02 & 15.06 & 27.09 & 25 & 26 & UNKNOWN\\
Qwen3.5 & IP-VCD & 27.13 & 71.43 & 48.19 & 42.75 & 80 & 32 & I cannot identify it\\
Qwen3.5 & IP-M3ID & 30.06 & 66.94 & 50.00 & 46.27 & 83 & 41 & UNKNOWN\\
Gemma-3 & Direct & 14.39 & 62.50 & 3.01 & 15.37 & 5 & 3 & UNKNOWN\\
Gemma-3 & VCD & 9.39 & 52.00 & 7.83 & 11.93 & 13 & 12 & NONE\\
Gemma-3 & DoLa & 8.96 & 84.62 & 6.63 & 11.11 & 11 & 2 & NONE\\
Gemma-3 & DeCo & 8.59 & 80.00 & 4.82 & 10.16 & 8 & 2 & NONE\\
Gemma-3 & SID & 12.64 & 66.67 & 4.82 & 14.20 & 8 & 4 & NONE\\
Gemma-3 & CDA & 4.75 & 53.17 & 40.36 & 17.83 & 67 & 59 & I cannot identify it\\
Gemma-3 & IP-VCD & 11.66 & 65.66 & 39.16 & 24.36 & 65 & 34 & UNCLEAR\\
Gemma-3 & IP-M3ID & 11.54 & 73.81 & 37.35 & 23.65 & 62 & 22 & UNKNOWN\\
OneVision & Direct & 30.61 & 100.00 & 7.23 & 32.95 & 12 & 0 & UNKNOWN\\
OneVision & VCD & 25.53 & 100.00 & 4.82 & 27.09 & 8 & 0 & NONE\\
OneVision & DoLa & 17.40 & 70.59 & 7.23 & 19.75 & 12 & 5 & NONE\\
OneVision & DeCo & 29.65 & 100.00 & 4.82 & 31.21 & 8 & 0 & NONE\\
OneVision & SID & 26.37 & 100.00 & 8.43 & 29.10 & 14 & 0 & NONE\\
OneVision & CDA & 26.43 & 79.49 & 18.67 & 32.48 & 31 & 8 & UNKNOWN\\
OneVision & IP-VCD & 25.84 & 76.70 & 47.59 & 41.27 & 79 & 24 & I cannot identify it\\
OneVision & IP-M3ID & 26.80 & 72.45 & 42.77 & 40.66 & 71 & 27 & UNSURE\\
Phi-3.5-Vision & Direct & 21.88 & 71.64 & 28.92 & 31.25 & 48 & 19 & UNKNOWN\\
Phi-3.5-Vision & VCD & 17.68 & 73.81 & 18.67 & 23.73 & 31 & 11 & NONE\\
Phi-3.5-Vision & DoLa & 18.26 & 67.14 & 28.31 & 27.44 & 47 & 23 & NONE\\
Phi-3.5-Vision & DeCo & 16.97 & 72.46 & 30.12 & 26.74 & 50 & 19 & NONE\\
Phi-3.5-Vision & SID & 17.89 & 78.95 & 18.07 & 23.75 & 30 & 8 & NONE\\
Phi-3.5-Vision & CDA & 17.70 & 63.38 & 54.22 & 35.27 & 90 & 52 & UNKNOWN\\
Phi-3.5-Vision & IP-VCD & 17.71 & 66.67 & 63.86 & 38.42 & 106 & 53 & I cannot identify it\\
Phi-3.5-Vision & IP-M3ID & 19.36 & 60.61 & 60.24 & 38.89 & 100 & 65 & I cannot identify it\\
\end{longtable}\endgroup

\section{控制实验与测量单位}\label{ux63a7ux5236ux5b9eux9a8cux4e0eux6d4bux91cfux5355ux4f4d}

\subsection{候选准入的九模型分解}\label{ux5019ux9009ux51c6ux5165ux7684ux4e5dux6a21ux578bux5206ux89e3}

代表面板按每模型每类一题构成909个首个生成位置，得到875个登记的答案／弃权首词元代理对。下表分别列出指定词元通过原生支持、引导支持、两候选均通过引导支持的数量，以及最后一组中参考指令交互压低弃权的数量。它们与完整回复的语义弃权是不同测量单位。

\begin{longtable}[]{@{}llllll@{}}
\toprule\noalign{}
模型 & 代理对 & 原生准入 & 引导准入 & 共同准入 & 参考压低 \\
\midrule\noalign{}
\endhead
\bottomrule\noalign{}
\endlastfoot
Gemma-3 & 88 & 0 & 0 & 0 & 0 \\
InternVL3.5 & 101 & 0 & 5 & 2 & 0 \\
LLaVA-Mistral & 94 & 0 & 67 & 57 & 44 \\
MiniCPM-V2.6 & 94 & 0 & 20 & 12 & 7 \\
OneVision & 101 & 0 & 1 & 1 & 0 \\
Phi-3.5-Vision & 101 & 0 & 6 & 3 & 2 \\
Qwen2.5-VL & 96 & 0 & 84 & 52 & 44 \\
Qwen3.5 & 99 & 0 & 60 & 56 & 37 \\
Qwen3-VL & 101 & 0 & 39 & 16 & 13 \\
\end{longtable}

\subsection{参考指令交叉控制}\label{ux53c2ux8003ux6307ux4ee4ux4ea4ux53c9ux63a7ux5236}

固定 Qwen2.5-VL 的正常图像及同一2,424题，行表示正常侧弃权表达，列表示参考侧表达，数值为语义弃权率（百分比）。

\begin{longtable}[]{@{}lllll@{}}
\toprule\noalign{}
正常与参考 & UNKNOWN & UNCLEAR & UNSURE & I cannot identify it \\
\midrule\noalign{}
\endhead
\bottomrule\noalign{}
\endlastfoot
UNKNOWN & 31.72 & 84.03 & 84.86 & 87.58 \\
UNCLEAR & 61.34 & 12.33 & 55.98 & 68.94 \\
UNSURE & 86.30 & 64.89 & 24.46 & 95.79 \\
I cannot identify it & 81.56 & 80.61 & 79.74 & 29.46 \\
\end{longtable}

\subsection{同指令与参考去引导比较}\label{ux540cux6307ux4ee4ux4e0eux53c2ux8003ux53bbux5f15ux5bfcux6bd4ux8f83}

九模型、四表达分别形成36个 VCD 与36个 M3ID 配对。匹配引导保留同一正常表达；参考去引导条件仅移除参考侧的弃权要求。完整72行配对 J 差值及216行条件指标随统计文件保存。VCD 的 IP 减双侧引导平均为0.4790个百分点，范围−2.7228至2.5165；IP 减参考去引导平均为3.1926，范围−8.8696至25.7838。M3ID 的相应平均为−0.0596与−1.3327个百分点，范围分别为−4.7442至3.6304和−7.5495至4.8680。

\subsection{匹配运行的计时范围}\label{ux5339ux914dux8fd0ux884cux7684ux8ba1ux65f6ux8303ux56f4}

计时比较筛选14个完整模型与表达匹配组，每组包含 VCD、IP-VCD 与 CDA 的2,424条记录。组内运行来源、GPU 身份、精度、软件版本与任务分配相同。时延包含图像读取和转换、提示与会话构造、生成及清理，排除模型加载及结果序列化；没有独立的 prefill/decode 计时或显存峰值记录。完整42行汇总同时保留平均生成长度、每生成词元墙钟时间、相对时延和软件设置。

\section{指令保持的数学性质与完整证明}\label{ux6307ux4ee4ux4fddux6301ux7684ux6570ux5b66ux6027ux8d28ux4e0eux5b8cux6574ux8bc1ux660e}

本附录从共享生成前缀下的条件分布出发，推导候选准入、参考指令交互与IP-VCD的保持性质，并给出三条件线性族的系数刻画、固定支持上的变分表达、候选组概率及完整回复概率的证明。VCD的视觉对比与候选约束采用Leng等人的定义\citep{leng2024vcd}；三条件构造及其输入职责沿用本文方法。下述等式均明确保留对应的候选支持与归一化。

\subsection{条件分布、候选支持与归一化}\label{ux6761ux4ef6ux5206ux5e03ux5019ux9009ux652fux6301ux4e0eux5f52ux4e00ux5316}

\subsubsection{四个条件与共享前缀}\label{ux56dbux4e2aux6761ux4ef6ux4e0eux5171ux4eabux524dux7f00}

设冻结的视觉语言模型为\(p_\theta\)，有限词表为\(\mathcal V\)且\(|\mathcal V|\geq2\)，正常图像为\(I\)，参考图像为\(\widetilde I\)，当前已生成的词元前缀为\(y_{<t}\)。在一次条件比较中，\(\widetilde I\)及其扰动随机量保持固定。令\(x_0\)包含任务问题与简洁回答要求，\(x_1\)在\(x_0\)上增加指定的弃权条件句。四个条件的下一词元对数概率定义为

\[
\begin{aligned}
c_t(v)&=\log p_\theta(v\mid I,x_0,y_{<t}),&
r_t(v)&=\log p_\theta(v\mid\widetilde I,x_0,y_{<t}),\\
g_t(v)&=\log p_\theta(v\mid I,x_1,y_{<t}),&
h_t(v)&=\log p_\theta(v\mid\widetilde I,x_1,y_{<t}).
\end{aligned}
\tag{D.1}
\]

四路使用相同模型、词表与实际回复前缀，各自保留与其输入条件对应的缓存。设各分支在\(\mathcal V\)上的未截断logit有限，因此其softmax概率为正。局部推导固定\(t\)并省略时间下标，记

\[
p_c=e^c,\quad p_r=e^r,\quad p_g=e^g,\quad p_h=e^h,
\qquad
\delta_c=g-c,\quad\delta_r=h-r,\quad d=c-r.
\tag{D.2}
\]

\(\delta_c\)与\(\delta_r\)逐词元记录同一弃权条件句在两种视觉输入上引起的对数概率变化。\(d\)比较正常图像与参考图像对同一续写的支持；它的两个条件都保留任务问题、简洁回答要求和当前回复前缀。

\subsubsection{候选约束与解码分布}\label{ux5019ux9009ux7ea6ux675fux4e0eux89e3ux7801ux5206ux5e03}

对对数概率向量\(\ell\)和\(0<\beta\leq1\)，定义

\[
\mathcal S_\ell
=\{v\in\mathcal V:\ell(v)\geq\max_{w\in\mathcal V}\ell(w)+\log\beta\}
=\{v:p_\ell(v)\geq\beta\max_w p_\ell(w)\}.
\tag{D.3}
\]

当\(\beta=0\)时，约定\(\mathcal S_\ell=\mathcal V\)。任意最大概率词元均属于\(\mathcal S_\ell\)，所以该集合非空。对非空集合\(S\subseteq\mathcal V\)和有限分数\(s\)，定义受限归一化算子

\[
\mathcal P_S[s](v)
=\frac{\mathbf1\{v\in S\}\exp s(v)}{Z_S(s)},
\qquad Z_S(s)=\sum_{w\in S}\exp s(w).
\tag{D.4}
\]

三个主要解码条件分别为

\[
\begin{aligned}
s_N&=(1+\alpha)c-\alpha r,&p_N&=\mathcal P_{\mathcal S_c}[s_N],\\
s_G&=(1+\alpha)g-\alpha h,&p_G&=\mathcal P_{\mathcal S_g}[s_G],\\
s_{\mathrm{IP}}&=g+\alpha(c-r),&p_{\mathrm{IP}}&=\mathcal P_{\mathcal S_g}[s_{\mathrm{IP}}],
\end{aligned}
\qquad \alpha\geq0.
\tag{D.5}
\]

下标\(N\)、\(G\)与\(\mathrm{IP}\)依次表示原生无弃权引导VCD、带引导VCD与IP-VCD。它们的分数函数定义于整个词表，实际归一化支持由式（D.5）指定。本文贪心生成在每一步选择该支持内分数最高的词元；并列时采用实现登记的固定顺序。后文的概率恒等式讨论式（D.4）所定义的归一化分数分布，词元选择由相应的最大值条件给出。

\textbf{引理D.1（归一化常数的消去）。} 对任意\(u,v\in S\)，有

\[
\log\frac{\mathcal P_S[s](u)}{\mathcal P_S[s](v)}=s(u)-s(v).
\tag{D.6}
\]

\textbf{证明。} 两个概率的分母同为\(Z_S(s)\)，取比值后分母抵消，再取对数即可。进一步，若第\(j\)个分支的logit为\(z_j\)，则\(\ell_j=z_j-\log\sum_v e^{z_j(v)}\)。对任意与词元无关的系数\(a_j\)，\(\sum_j a_jz_j\)与\(\sum_j a_j\ell_j\)相差一个与词元无关的常数。受限softmax、候选间的分数差及argmax均保持相同。候选准入比较\(\ell_j(v)-\max_w\ell_j(w)\)，同样消去该常数。证毕。

这一性质使原始logit实现与对数概率推导对应起来，也使以下保持性质可以写成实际归一化分布的对数优势比。

\subsection{两阶段作用：候选准入与参考指令交互}\label{ux4e24ux9636ux6bb5ux4f5cux7528ux5019ux9009ux51c6ux5165ux4e0eux53c2ux8003ux6307ux4ee4ux4ea4ux4e92}

\subsubsection{弃权候选的准入条件}\label{ux5f03ux6743ux5019ux9009ux7684ux51c6ux5165ux6761ux4ef6}

对\(0<\beta\leq1\)，定义候选\(u\)距准入阈值的差

\[
K_\ell(u)=\ell(u)-\max_w\ell(w)-\log\beta.
\tag{D.7}
\]

\(K_\ell(u)\geq0\)恰好对应\(u\in\mathcal S_\ell\)。由于\(g=c+\delta_c\)，有

\[
K_g(u)-K_c(u)
=\delta_c(u)-\big[\max_w\{c(w)+\delta_c(w)\}-\max_w c(w)\big].
\tag{D.8}
\]

\textbf{证明。} 将\(g=c+\delta_c\)分别代入\(K_g\)与\(K_c\)，消去\(c(u)\)和\(\log\beta\)即得。等价地，候选准入可以写成所有竞争词元上的条件

\[
u\in\mathcal S_g
\iff
[c(u)-c(v)]+[\delta_c(u)-\delta_c(v)]\geq\log\beta,
\quad\forall v\in\mathcal V.
\tag{D.9}
\]

因为\(g(u)\geq\max_vg(v)+\log\beta\)等价于\(g(u)-g(v)\geq\log\beta\)对所有\(v\)成立。证毕。

式（D.8）表明，指令对弃权候选的增量要与词表最大分数的移动共同比较。候选的相对位置越过阈值后，才进入后续的分数竞争。由式（D.4），支持之外候选的概率为零；支持内候选的相对概率由\(s\)确定。IP-VCD的准入集合由\(g\)给出，视觉参考\(r\)参与其上的分数调整。

\subsubsection{带引导VCD的完整分解}\label{ux5e26ux5f15ux5bfcvcdux7684ux5b8cux6574ux5206ux89e3}

\textbf{命题D.1（参考指令交互）。} 四路条件满足

\[
\boxed{
s_G=s_N+\delta_c+\alpha(\delta_c-\delta_r),\qquad
s_{\mathrm{IP}}=s_N+\delta_c.
}
\tag{D.10}
\]

\textbf{证明。} 使用\(g=c+\delta_c\)和\(h=r+\delta_r\)，逐项展开：

\[
\begin{aligned}
s_G
&=(1+\alpha)(c+\delta_c)-\alpha(r+\delta_r)\\
&=(1+\alpha)c-\alpha r+(1+\alpha)\delta_c-\alpha\delta_r\\
&=s_N+\delta_c+\alpha(\delta_c-\delta_r),\\
s_{\mathrm{IP}}
&=c+\delta_c+\alpha(c-r)
=s_N+\delta_c.
\end{aligned}
\tag{D.11}
\]

两式相减得到

\[
s_{\mathrm{IP}}-s_G=\alpha(\delta_r-\delta_c).
\tag{D.12}
\]

证毕。该分解将原生视觉对比、正常图像上的一次指令增量以及额外的参考交互分别列出。相同指令在不同视觉条件上的作用由两条实际条件分布给定。

\subsubsection{从分数变化到候选翻转}\label{ux4eceux5206ux6570ux53d8ux5316ux5230ux5019ux9009ux7ffbux8f6c}

固定一个弃权续写候选\(u\)和一个答案续写候选\(v\)，统一以``弃权减答案''为边际方向：

\[
m_\ell=\ell(u)-\ell(v),\qquad
\Delta_c=\delta_c(u)-\delta_c(v),\qquad
\Delta_r=\delta_r(u)-\delta_r(v).
\tag{D.13}
\]

对组合分数同样记\(m_N=s_N(u)-s_N(v)\)等。于是

\[
\begin{aligned}
m_N&=m_c+\alpha(m_c-m_r),\\
m_G&=m_N+\Delta_c+\alpha(\Delta_c-\Delta_r),\\
m_{\mathrm{IP}}&=m_N+\Delta_c
=m_G+\alpha(\Delta_r-\Delta_c).
\end{aligned}
\tag{D.14}
\]

当\(\alpha>0\)且\(\Delta_r>\Delta_c\)时，参考条件中的额外指令响应使带引导VCD的弃权边际减少\(\alpha(\Delta_r-\Delta_c)\)。当\(u,v\in\mathcal S_g\)时，由引理D.1可得归一化概率的精确关系

\[
\frac{p_{\mathrm{IP}}(u)/p_{\mathrm{IP}}(v)}{p_G(u)/p_G(v)}
=\exp\{\alpha(\Delta_r-\Delta_c)\}.
\tag{D.15}
\]

因此，两候选之间由答案领先变为弃权领先的严格条件为

\[
m_G<0<m_G+\alpha(\Delta_r-\Delta_c).
\tag{D.16}
\]

若讨论整个候选集合中的贪心胜出者，令\(k(w)=\alpha[\delta_r(w)-\delta_c(w)]\)。候选\(u\)成为IP-VCD的唯一最大分数词元，当且仅当

\[
s_G(u)-s_G(v)+k(u)-k(v)>0,
\qquad\forall v\in\mathcal S_g\setminus\{u\}.
\tag{D.17}
\]

式（D.16）描述指定候选对的翻转，式（D.17）描述全词表保留集合中的胜出条件。两者分别对应局部边际记录与实际argmax记录。

无引导参考也可以改变弃权的相对支持：由式（D.14），\(m_r>m_c\)且\(\alpha>0\)时，原生对比相对正常分布的边际改变量为\(\alpha(m_c-m_r)<0\)。这项视觉差分继续包含在IP-VCD中；式（D.12）精确列出IP-VCD对额外指令交互的处理。

\subsubsection{仅移除参考引导的对照}\label{ux4ec5ux79fbux9664ux53c2ux8003ux5f15ux5bfcux7684ux5bf9ux7167}

仅在参考侧移除弃权条件句，其分数为

\[
\begin{aligned}
s_{\mathrm{ref\text{-}off}}
&=(1+\alpha)g-\alpha r\\
&=s_N+(1+\alpha)\delta_c
=s_{\mathrm{IP}}+\alpha\delta_c.
\end{aligned}
\tag{D.18}
\]

该对照保留了\(1+\alpha\)倍的正常指令增量。IP-VCD增加的无引导正常条件\(c\)，使视觉差分的系数\(\alpha\)与正常指令的单位系数分别确定。

\subsection{指令保持性质与三条件系数}\label{ux6307ux4ee4ux4fddux6301ux6027ux8d28ux4e0eux4e09ux6761ux4ef6ux7cfbux6570}

\subsubsection{相对指令作用的精确保持}\label{ux76f8ux5bf9ux6307ux4ee4ux4f5cux7528ux7684ux7cbeux786eux4fddux6301}

\textbf{定理D.1（单位指令增量保持）。} 固定生成前缀，并在同一个非空支持\(S\)上分别归一化\(s_N\)和\(s_{\mathrm{IP}}\)。对任意\(u,v\in S\)，有

\[
\boxed{
\log\frac{\mathcal P_S[s_{\mathrm{IP}}](u)}{\mathcal P_S[s_{\mathrm{IP}}](v)}
-
\log\frac{\mathcal P_S[s_N](u)}{\mathcal P_S[s_N](v)}
=\delta_c(u)-\delta_c(v).
}
\tag{D.19}
\]

\textbf{证明。} 由引理D.1，左侧等于

\[
\begin{aligned}
&[s_{\mathrm{IP}}(u)-s_{\mathrm{IP}}(v)]-[s_N(u)-s_N(v)]\\
&=[s_N(u)+\delta_c(u)-s_N(v)-\delta_c(v)]-[s_N(u)-s_N(v)]\\
&=\delta_c(u)-\delta_c(v).
\end{aligned}
\]

证毕。对于各自采用\(\mathcal S_c\)和\(\mathcal S_g\)的实际分布，同一结论适用于\(u,v\in\mathcal S_c\cap\mathcal S_g\)。对新进入\(\mathcal S_g\)的候选，采用\(s_N\)在\(\mathcal S_g\)上的受限归一化，可以单独观察同一候选集合中的指令作用。

将式（D.19）指数化，有

\[
\frac{\mathcal P_S[s_{\mathrm{IP}}](u)}{\mathcal P_S[s_{\mathrm{IP}}](v)}
=
\frac{\mathcal P_S[s_N](u)}{\mathcal P_S[s_N](v)}
\frac{p_g(u)/p_c(u)}{p_g(v)/p_c(v)}.
\tag{D.20}
\]

因此，正常图像上由指令引起的候选优势比，完整乘入原生视觉对比的候选优势比中。对应地，相对带引导正常基础\(p_g\)，IP-VCD施加的视觉增量为

\[
\log\frac{p_{\mathrm{IP}}(u)}{p_{\mathrm{IP}}(v)}
-\log\frac{p_g(u)}{p_g(v)}
=\alpha[d(u)-d(v)],\qquad u,v\in\mathcal S_g.
\tag{D.21}
\]

式（D.19）与式（D.21）共同表达两项职责：正常指令以单位系数进入候选竞争，视觉差分以\(\alpha\)为强度进行调整。

\subsubsection{三条件线性组合中的唯一系数}\label{ux4e09ux6761ux4ef6ux7ebfux6027ux7ec4ux5408ux4e2dux7684ux552fux4e00ux7cfbux6570}

\textbf{定理D.2（三条件线性族的系数刻画）。} 考虑与词元无关的系数组成的三条件线性族

\[
s_{a,b,e}=a g+b c+e r.
\tag{D.22}
\]

要求该族对所有正条件分布满足以下两项性质。其一，固定\(c,r\)而将\(g\)替换为\(g'\)时，任意候选对的分数改变量等于该候选对在\(g'-g\)中的改变量。其二，在\(g=c\)时，候选对分数差恢复为原生VCD的分数差。则系数唯一为

\[
(a,b,e)=(1,\alpha,-\alpha).
\tag{D.23}
\]

\textbf{证明。} 第一项要求给出

\[
a\big[(g'(u)-g(u))-(g'(v)-g(v))\big]
=(g'(u)-g(u))-(g'(v)-g(v)).
\]

正分布的候选对对数优势比可以连续变化，因此该恒等式要求\(a=1\)。第二项要求在\(g=c\)时，对任意\(c,r,u,v\)满足

\[
(a+b)[c(u)-c(v)]+e[r(u)-r(v)]
=(1+\alpha)[c(u)-c(v)]-\alpha[r(u)-r(v)].
\]

分别变化\(c\)与\(r\)的候选对优势比，得到\(a+b=1+\alpha\)及\(e=-\alpha\)。代入\(a=1\)，即得\(b=\alpha\)。证毕。

唯一性针对式（D.22）的线性族及上述两项设计要求。所有分数共同增加一个与词元无关的常数时，解码分布保持相同。

\subsubsection{概率形式与退化情形}\label{ux6982ux7387ux5f62ux5f0fux4e0eux9000ux5316ux60c5ux5f62}

将三条件分数代入受限归一化，得到

\[
p_{\mathrm{IP}}(v)
=\frac{\mathbf1\{v\in\mathcal S_g\}\,p_g(v)
\left[p_c(v)/p_r(v)\right]^\alpha}
{\sum_{w\in\mathcal S_g}p_g(w)\left[p_c(w)/p_r(w)\right]^\alpha}.
\tag{D.24}
\]

这种基础分布与似然比相乘的形式属于解码期专家乘积构造\citep{liu2021dexperts}。这里的三个分布均由同一冻结模型在已明确的三个输入条件下产生。

若\(g=c\)，则\(\mathcal S_g=\mathcal S_c\)，分数和支持共同恢复原生VCD。若\(c=r\)或\(\alpha=0\)，则

\[
p_{\mathrm{IP}}(v)
=\frac{\mathbf1\{v\in\mathcal S_g\}p_g(v)}{p_g(\mathcal S_g)}.
\tag{D.25}
\]

此时得到引导基础在保留集合上的条件分布。由于\(\mathcal S_g\)包含全部最大概率候选，贪心选择与引导Direct相同。若\(\beta=0\)，式（D.25）进一步等于完整的\(p_g\)。

\subsection{KL正则化下的唯一最优分布}\label{klux6b63ux5219ux5316ux4e0bux7684ux552fux4e00ux6700ux4f18ux5206ux5e03}

\subsubsection{目标函数与闭式解}\label{ux76eeux6807ux51fdux6570ux4e0eux95edux5f0fux89e3}

固定生成前缀及支持\(S=\mathcal S_g\)，记引导基础在该支持上的条件分布为

\[
G(v)=\frac{p_g(v)}{p_g(S)},\qquad v\in S,
\tag{D.26}
\]

并令\(\Delta(S)\)表示支持\(S\)上的概率单纯形。视觉相对支持为\(d(v)=c(v)-r(v)\)。考虑局部目标

\[
\mathcal L_\alpha(\pi)
=\alpha\sum_{v\in S}\pi(v)d(v)
-\sum_{v\in S}\pi(v)\log\frac{\pi(v)}{G(v)},
\qquad\pi\in\Delta(S),
\tag{D.27}
\]

其中约定\(0\log0=0\)。第一项奖励正常视觉相对参考视觉的支持，第二项以KL散度约束相对引导基础的变化。

\textbf{定理D.3（固定支持上的变分最优性）。} 目标（D.27）的唯一最大化分布为

\[
\pi_\alpha^*(v)
=\frac{G(v)e^{\alpha d(v)}}{\overline Z_\alpha},
\qquad
\overline Z_\alpha=\sum_{w\in S}G(w)e^{\alpha d(w)}.
\tag{D.28}
\]

该分布恰为\(p_{\mathrm{IP}}\)。

\textbf{证明。} 由式（D.28），

\[
\log\pi_\alpha^*(v)=\log G(v)+\alpha d(v)-\log\overline Z_\alpha.
\]

对任意\(\pi\in\Delta(S)\)，代入KL散度：

\[
\begin{aligned}
D_{\mathrm{KL}}(\pi\|\pi_\alpha^*)
&=\sum_v\pi(v)[\log\pi(v)-\log G(v)-\alpha d(v)+\log\overline Z_\alpha]\\
&=D_{\mathrm{KL}}(\pi\|G)-\alpha\mathbb E_\pi[d]+\log\overline Z_\alpha.
\end{aligned}
\]

整理得

\[
\boxed{
\mathcal L_\alpha(\pi)
=\log\overline Z_\alpha-D_{\mathrm{KL}}(\pi\|\pi_\alpha^*).
}
\tag{D.29}
\]

KL非负性可直接由\(\log x\leq x-1\)推出。对\(\pi(v)>0\)的项求和，有

\[
-\sum_{v:\pi(v)>0}\pi(v)\log\frac{\pi_\alpha^*(v)}{\pi(v)}
\geq\sum_{v:\pi(v)>0}[\pi(v)-\pi_\alpha^*(v)]\geq0.
\]

等号要求全部正概率位置的比值为1，且\(\pi_\alpha^*\)在其余位置的总质量为0。由于\(\pi_\alpha^*\)在\(S\)上处处为正，等号恰在\(\pi=\pi_\alpha^*\)时成立。因此（D.29）的唯一最大值为\(\log\overline Z_\alpha\)，由\(\pi_\alpha^*\)取得。再将\(G=p_g/p_g(S)\)代入（D.28），公共因子\(p_g(S)\)抵消，得到（D.24）。证毕。

这一最优化解释将IP-VCD对应到固定前缀上的明确目标。用于实验的联合决策效用采用完整回答及冻结参考标签计算，两者分别承担局部构造与最终行为评价的职责。

\subsubsection{对比强度怎样改变分布}\label{ux5bf9ux6bd4ux5f3aux5ea6ux600eux6837ux6539ux53d8ux5206ux5e03}

令\(\psi(\alpha)=\log\overline Z_\alpha\)。有限词表允许逐项求导，得到

\[
\psi'(\alpha)=\mathbb E_{\pi_\alpha^*}[d],\qquad
\psi''(\alpha)=\operatorname{Var}_{\pi_\alpha^*}(d)\geq0.
\tag{D.30}
\]

\textbf{证明。} 对\(\overline Z_\alpha\)求导并除以自身，得到第一式。继续求导，第二式为\(\mathbb E[d^2]-(\mathbb E[d])^2\)。等价地，对每个候选有

\[
\frac{\partial\pi_\alpha^*(v)}{\partial\alpha}
=\pi_\alpha^*(v)\big[d(v)-\mathbb E_{\pi_\alpha^*}[d]\big].
\tag{D.31}
\]

证毕。因此，在固定四路分布与固定支持下，增大\(\alpha\)提高视觉相对支持高于当前均值的候选概率，并使其期望单调不减。

对预先定义的候选组\(A\subset S\)，令\(B=S\setminus A\)，且两组非空。求和得到

\[
\frac{\partial\pi_\alpha^*(A)}{\partial\alpha}
=\pi_\alpha^*(A)\pi_\alpha^*(B)
\left[\mathbb E_{\pi_\alpha^*}[d\mid A]-\mathbb E_{\pi_\alpha^*}[d\mid B]\right].
\tag{D.32}
\]

候选组的概率增长由两组的条件平均视觉支持之差决定，这给出了对比强度作用于整组候选的精确条件。

\subsection{从单个候选到候选组概率}\label{ux4eceux5355ux4e2aux5019ux9009ux5230ux5019ux9009ux7ec4ux6982ux7387}

\subsubsection{指数倾斜的一般分组恒等式}\label{ux6307ux6570ux503eux659cux7684ux4e00ux822cux5206ux7ec4ux6052ux7b49ux5f0f}

令\(p\)是\(\mathcal V\)上的正分布，\(A,B\)构成\(\mathcal V\)的非空划分，且\(A\cap S\)、\(B\cap S\)均非空。对任意有限函数\(f\)，定义

\[
q(v)=\frac{\mathbf1\{v\in S\}p(v)e^{f(v)}}{\sum_{w\in S}p(w)e^{f(w)}}.
\tag{D.33}
\]

记\(O_p(A)=p(A)/p(B)\)，并令\(p_S=p(\cdot\mid S)\)。有

\[
\boxed{
\begin{aligned}
\log O_q(A)-\log O_p(A)
={}&\log\frac{p(S\mid A)}{p(S\mid B)}\\
&+\log\mathbb E_{p_S(\cdot\mid A)}[e^f]
-\log\mathbb E_{p_S(\cdot\mid B)}[e^f].
\end{aligned}}
\tag{D.34}
\]

\textbf{证明。} 对两组分别求和，归一化常数相消：

\[
\begin{aligned}
\frac{q(A)}{q(B)}
&=\frac{\sum_{v\in A\cap S}p(v)e^{f(v)}}{\sum_{v\in B\cap S}p(v)e^{f(v)}}\\
&=\frac{p(A)p(S\mid A)}{p(B)p(S\mid B)}
\frac{\mathbb E_{p(\cdot\mid A\cap S)}[e^f]}{\mathbb E_{p(\cdot\mid B\cap S)}[e^f]}.
\end{aligned}
\]

取对数即得（D.34）。证毕。第一项刻画候选截断改变了多少组间优势比，其余两项刻画保留候选上的重新加权。若\(A\cap S\)为空，则\(q(A)=0\)；若\(B\cap S\)为空，则\(q(A)=1\)，相应概率可直接从支持确定。

对IP-VCD，取\(p=p_g\)、\(f=\alpha(c-r)\)及\(S=\mathcal S_g\)，式（D.34）即给出任意登记弃权词元组相对引导正常分布的完整概率变化。

\subsubsection{移除参考交互后的组概率变化}\label{ux79fbux9664ux53c2ux8003ux4ea4ux4e92ux540eux7684ux7ec4ux6982ux7387ux53d8ux5316}

带引导VCD与IP-VCD共享\(\mathcal S_g\)。令\(k=\alpha(\delta_r-\delta_c)\)，由式（D.12）有

\[
p_{\mathrm{IP}}(v)
=\frac{p_G(v)e^{k(v)}}{\mathbb E_{p_G}[e^k]},\qquad v\in\mathcal S_g.
\tag{D.35}
\]

因此，对保留集合中的两组\(A,B\)，

\[
\log O_{p_{\mathrm{IP}}}(A)-\log O_{p_G}(A)
=\log\mathbb E_{p_G(\cdot\mid A)}[e^k]
-\log\mathbb E_{p_G(\cdot\mid B)}[e^k].
\tag{D.36}
\]

\textbf{证明。} 将\(s_{\mathrm{IP}}=s_G+k\)代入受限softmax，分母除以\(Z_{\mathcal S_g}(s_G)\)，得到（D.35）；随后使用（D.34）在同一支持上的形式即可。证毕。

组概率严格增加的充要条件为

\[
p_{\mathrm{IP}}(A)>p_G(A)
\iff
\mathbb E_{p_G(\cdot\mid A)}[e^k]
>\mathbb E_{p_G(\cdot\mid B)}[e^k].
\tag{D.37}
\]

一个直观的充分条件是\(\min_{u\in A}k(u)>\max_{v\in B}k(v)\)。对于单个候选，精确条件为\(k(u)>\log\mathbb E_{p_G}[e^k]\)。这些关系将局部加分与归一化后的概率变化连接起来：提高组间优势需要比较两组候选的加权增量。

\subsubsection{VCD的候选支持、组间偏好与组内分配分解}\label{vcdux7684ux5019ux9009ux652fux6301ux7ec4ux95f4ux504fux597dux4e0eux7ec4ux5185ux5206ux914dux5206ux89e3}

令\(p,q\)分别为某一VCD实例的正常与参考概率，支持为\(S\)，对比分布为

\[
m(v)\propto\mathbf1\{v\in S\}p(v)^{1+\alpha}q(v)^{-\alpha}.
\tag{D.38}
\]

定义\(p_S=p(\cdot\mid S)\)、\(q_S=q(\cdot\mid S)\)，并在两个保留组上定义\(p_A=p_S(\cdot\mid A)\)、\(q_A=q_S(\cdot\mid A)\)及\(p_B,q_B\)。对\(\rho>1\)，Rényi散度采用\citep{vanerven2014renyi}的有限分布定义

\[
D_\rho(P\|Q)=\frac{1}{\rho-1}\log\sum_v P(v)^\rho Q(v)^{1-\rho}.
\tag{D.39}
\]

\textbf{命题D.2（整组对比作用的精确分解）。} 当\(\alpha>0\)且两组在\(S\)上都有正质量时，

\[
\boxed{
\begin{aligned}
\log O_m(A)-\log O_p(A)
={}&\underbrace{\log O_{p_S}(A)-\log O_p(A)}_{\text{候选支持作用}}\\
&+\underbrace{\alpha[\log O_{p_S}(A)-\log O_{q_S}(A)]}_{\text{参考组间偏好作用}}\\
&+\underbrace{\alpha[D_{1+\alpha}(p_A\|q_A)-D_{1+\alpha}(p_B\|q_B)]}_{\text{组内候选分配作用}}.
\end{aligned}}
\tag{D.40}
\]

\textbf{证明。} 将\(p,q\)替换为\(p_S,q_S\)只引入一个组间相同的比例常数。组\(A\)上的未归一化质量可写为

\[
\begin{aligned}
F_A
&=\sum_{v\in A\cap S}p_S(v)^{1+\alpha}q_S(v)^{-\alpha}\\
&=p_S(A)^{1+\alpha}q_S(A)^{-\alpha}
\sum_{v\in A\cap S}p_A(v)^{1+\alpha}q_A(v)^{-\alpha}\\
&=p_S(A)^{1+\alpha}q_S(A)^{-\alpha}
\exp\{\alpha D_{1+\alpha}(p_A\|q_A)\}.
\end{aligned}
\tag{D.41}
\]

对\(B\)同理。由\(O_m(A)=F_A/F_B\)，得

\[
\log O_m(A)
=(1+\alpha)\log O_{p_S}(A)-\alpha\log O_{q_S}(A)
+\alpha[D_{1+\alpha}(p_A\|q_A)-D_{1+\alpha}(p_B\|q_B)].
\]

减去\(\log O_p(A)\)，再拆分\((1+\alpha)\log O_{p_S}(A)\)即得（D.40）。证毕。当\(\alpha=0\)时，\(m=p_S\)，概率变化由候选支持项给出。

该分解将整组弃权变化落实到三个可计算对象：哪些候选通过阈值、参考分布怎样偏向两个组，以及各组内部怎样分配不同表达的概率。当保留集合只含两个候选、每组各一个时，两个组内散度均为零，式（D.40）恢复候选对的边际分解。局部弃权组由登记词元集合确定；完整回复的语义弃权使用下一节的终止序列事件定义。

\subsection{自回归生成中的完整回复概率}\label{ux81eaux56deux5f52ux751fux6210ux4e2dux7684ux5b8cux6574ux56deux590dux6982ux7387}

\subsubsection{局部归一化沿路径累积}\label{ux5c40ux90e8ux5f52ux4e00ux5316ux6cbfux8defux5f84ux7d2fux79ef}

设\(\mathcal Y_T\)为在首个登记终止词元处结束、或到达最大长度\(T\)时结束的所有输出词元序列。对任意方法\(M\)，其逐步归一化分数诱导

\[
P_M(y)=\prod_{t=1}^{|y|}p_{M,t}(y_t\mid y_{<t}),\qquad y\in\mathcal Y_T.
\tag{D.42}
\]

生成树的每个内部结点将其概率质量按归一化的下一词元概率分配给子结点，因此所有叶结点的概率和为1。对于由终止词元结束的回复，乘积包含该终止词元；对于达到长度上限的回复，乘积包含实际生成的最后一个词元。

对IP-VCD，定义

\[
Z_{\mathrm{IP},t}(y_{<t})
=\sum_{v\in\mathcal S_{g,t}(y_{<t})}
\exp\{g_t(v)+\alpha[c_t(v)-r_t(v)]\}.
\tag{D.43}
\]

对每一步均通过支持约束的路径，代入（D.42）并取对数：

\[
\begin{aligned}
\log P_{\mathrm{IP}}(y)
={}&\sum_t g_t(y_t)
+\alpha\sum_t[c_t(y_t)-r_t(y_t)]\\
&-\sum_t\log Z_{\mathrm{IP},t}(y_{<t}).
\end{aligned}
\tag{D.44}
\]

路径一旦包含支持外词元，其\(P_{\mathrm{IP}}(y)=0\)。各条件分布在式（D.44）中均沿同一条\(y\)计算，局部归一化常数也由该路径的前缀确定。

进一步令\(P_g(y)=\prod_t e^{g_t(y_t)}\)、\(P_c(y)=\prod_t e^{c_t(y_t)}\)、\(P_r(y)=\prod_t e^{r_t(y_t)}\)，并令\(\chi_g(y)=\prod_t\mathbf1\{y_t\in\mathcal S_{g,t}(y_{<t})\}\)。固定\(\alpha\)时，完整乘积形式为

\[
P_{\mathrm{IP}}(y)
=\frac{\chi_g(y)P_g(y)[P_c(y)/P_r(y)]^\alpha}
{\prod_t Z_{\mathrm{IP},t}(y_{<t})}.
\tag{D.45}
\]

式（D.45）保留了每个前缀对应的分母。实际回复重放逐步记录这些量，便能把局部分数调整累计为完整序列概率。

\subsubsection{IP-VCD与带引导VCD的路径差异}\label{ip-vcdux4e0eux5e26ux5f15ux5bfcvcdux7684ux8defux5f84ux5deeux5f02}

记带引导VCD对应的逐步归一化常数为\(Z_{G,t}\)。两个方法在同一给定前缀上使用相同的\(\mathcal S_{g,t}\)。对共同正概率路径，由（D.12）与（D.44）有

\[
\begin{aligned}
\log\frac{P_{\mathrm{IP}}(y)}{P_G(y)}
={}&\alpha\sum_t[\delta_{r,t}(y_t)-\delta_{c,t}(y_t)]\\
&-\sum_t\log\frac{Z_{\mathrm{IP},t}(y_{<t})}{Z_{G,t}(y_{<t})}.
\end{aligned}
\tag{D.46}
\]

\textbf{证明。} 对每一步写出\(\log p_{\mathrm{IP},t}-\log p_{G,t}\)，利用分数差（D.12），再沿路径求和。证毕。

式（D.46）将参考交互的移除量与每一步重新分配概率质量的作用同时计入。贪心路径从首次分岔后形成各自的前缀，完整路径比较通过逐条教师强制重放计算上述概率。

\subsubsection{完整弃权事件的概率变化}\label{ux5b8cux6574ux5f03ux6743ux4e8bux4ef6ux7684ux6982ux7387ux53d8ux5316}

令\(\mathcal A\subset\mathcal Y_T\)为按完整回复语义标记的弃权事件，\(\mathcal B=\mathcal Y_T\setminus\mathcal A\)。取在\(\mathcal Y_T\)上为正的基础路径分布\(P_0\)，例如未截断的引导条件分布\(P_g\)；取另一归一化路径分布\(P_1\)。定义路径的对数似然变化

\[
D(y)=\log\frac{P_1(y)}{P_0(y)},
\qquad e^{D(y)}=0\ \text{当}\ P_1(y)=0.
\tag{D.47}
\]

只要两个事件在所比较的分布上均有正质量，就有

\[
\boxed{
\log\frac{P_1(\mathcal A)}{P_1(\mathcal B)}
-\log\frac{P_0(\mathcal A)}{P_0(\mathcal B)}
=\log\mathbb E_{P_0(\cdot\mid\mathcal A)}[e^D]
-\log\mathbb E_{P_0(\cdot\mid\mathcal B)}[e^D].
}
\tag{D.48}
\]

\textbf{证明。} 由\(P_1(y)=P_0(y)e^{D(y)}\)，对事件\(\mathcal A\)求和得到

\[
P_1(\mathcal A)=P_0(\mathcal A)\mathbb E_{P_0(\cdot\mid\mathcal A)}[e^D].
\]

对\(\mathcal B\)同理，两式取比值并取对数即得。证毕。

若实际重放使用登记的有限完整回复集合\(\mathcal C\subseteq\mathcal Y_T\)，分别构造\(P_j(\cdot\mid\mathcal C)\)，并以\(\mathcal A\cap\mathcal C\)和\(\mathcal B\cap\mathcal C\)分组，则（D.48）同样成立。每条候选记录保存实际词元序列、结束方式及其完整语义标签，相同词元路径按一次计数。这样，单个词元的准入、指定候选对的翻转和完整弃权事件分别在各自的样本空间中具有明确的计算表达。

\subsection{向IP-M3ID推广}\label{ux5411ip-m3idux63a8ux5e7f}

M3ID使用带视觉条件与去图像条件之间的差分，并通过下一词元置信度和时间权重控制差分强度\citep{favero2024m3id}。在IP-M3ID中，\(c_t\)继续表示无弃权引导的正常视觉对数概率，\(r_t\)改为保留相同文本及生成前缀的去图像对数概率。令\(\tau_t\)为实现登记的时间索引，置信度阈值为\(\kappa\)，则当前构造的权重可写为

\[
w_t^c
=\mathbf1\{\max_v e^{c_t(v)}<\kappa\}
\big(e^{\lambda\tau_t}-1\big),
\qquad
\lambda\geq0,\quad\tau_t\geq0.
\tag{D.49}
\]

本研究按登记位置计算权重：若生成循环采用零起始索引\(j=t-1\)，则\(\tau_t=j+k_0\)，\(k_0\)为当前输入无弃权引导任务文本的登记词元长度，因此可随问题改变。采用\(\lambda=0.02\)、\(\kappa=0.3\)。权重由\(c_t\)及该索引决定。在固定前缀上，它对所有候选相同。

IP-M3ID及其同权重无引导形式分别为

\[
\begin{aligned}
s_{\mathrm{IPM},t}&=g_t+w_t^c(c_t-r_t),\\
s_{\mathrm{NM},t}&=c_t+w_t^c(c_t-r_t),\\
s_{\mathrm{IPM},t}-s_{\mathrm{NM},t}&=\delta_{c,t}.
\end{aligned}
\tag{D.50}
\]

本文的IP-M3ID分布定义在完整词表上归一化（等价于不施加VCD候选截断）。于是定理D.1的候选对保持性质对全词表成立；定理D.3在每个固定前缀上以\(w_t^c\)替换\(\alpha\)成立。路径概率则按每一步实际权重保留\(\sum_t w_t^c[c_t(y_t)-r_t(y_t)]\)。

对带引导M3ID，令\(w_t^g=\mathbf1\{\max_v e^{g_t(v)}<\kappa\}(e^{\lambda\tau_t^g}-1)\)，其中\(\tau_t^g\)为该引导条件的登记位置；以带引导去图像分布\(h_t\)作参考，则

\[
s_{\mathrm{GM},t}=g_t+w_t^g(g_t-h_t).
\]

利用\(g_t-h_t=(c_t-r_t)+(\delta_{c,t}-\delta_{r,t})\)，可得完整差值

\[
\boxed{
s_{\mathrm{IPM},t}-s_{\mathrm{GM},t}
=(w_t^c-w_t^g)(c_t-r_t)
+w_t^g(\delta_{r,t}-\delta_{c,t}).
}
\tag{D.51}
\]

第一项来自实际门控与时间权重的差异，第二项来自参考指令交互。取\(w_t^c=w_t^g\)时，公式恢复同权重的指令交互分解。该形式允许在保持真实调度规则的同时，分别核算两项来源。

\subsection{生牛肉片案例的逐项代入}\label{ux751fux725bux8089ux7247ux6848ux4f8bux7684ux9010ux9879ux4ee3ux5165}

取正文LLaVA-Mistral的生牛肉片案例，在首个生成位置比较弃权词元``UN''与答案词元``To''。保持实际输入、前缀和噪声不变，已有四路记录给出以下弃权减答案边际。

\begin{longtable}[]{@{}lrrrr@{}}
\toprule\noalign{}
条件 & \(m_c\) & \(m_r\) & \(m_g\) & \(m_h\) \\
\midrule\noalign{}
\endhead
\bottomrule\noalign{}
\endlastfoot
对数概率边际 & \(-12.5625\) & \(-11.546875\) & \(1.4375\) & \(3.0625\) \\
\end{longtable}

由式（D.13），

\[
\begin{aligned}
\Delta_c&=m_g-m_c=1.4375-(-12.5625)=14,\\
\Delta_r&=m_h-m_r=3.0625-(-11.546875)=14.609375.
\end{aligned}
\tag{D.52}
\]

在\(\alpha=1\)下，原生视觉对比的边际为

\[
m_N=2m_c-m_r=-13.578125.
\tag{D.53}
\]

带引导VCD与IP-VCD分别给出

\[
\begin{aligned}
m_G&=m_N+\Delta_c+(\Delta_c-\Delta_r)\\
&=-13.578125+14-0.609375=-0.1875,\\
m_{\mathrm{IP}}&=m_N+\Delta_c=-13.578125+14=0.421875.
\end{aligned}
\tag{D.54}
\]

两个候选均在引导支持内，因而其优势比变化为\(\exp(0.609375)\)。该位置的全词表最大分数词元也从``To''变为``UN''，对应真实完整回复``Toast''与``UNKNOWN''。

进一步在共同支持上插入诊断参数\(\eta\in[0,1]\)，定义

\[
s_\eta=s_{\mathrm{IP}}+\eta\alpha(\delta_c-\delta_r),
\qquad
m_\eta=0.421875-0.609375\eta.
\tag{D.55}
\]

固定前缀、共同支持上候选对分数的精确交叉点为

\[
\eta_* = \frac{0.421875}{0.609375}=\frac9{13}\approx0.692308.
\tag{D.56}
\]

\(\eta<9/13\)时，该候选对中弃权领先；\(\eta>9/13\)时，答案领先。已有的\(\eta=0,0.5,1\)诊断分别选择``UN''``UN''``To''。这一实例将正常指令增量、参考额外响应、共同支持上的分数翻转与实际回复路径连接起来。

\IfFileExists{appendix_supplement_v3.tex}{\section{开发集表达选择与补充机制结果}
\label{sec:supp-development-mechanism}

\subsection{开发集选择的实际工作点}
\label{sec:supp-dev-selected}

开发集选择在五个模型上使用每类4题、共404题，并为IP-VCD、同措辞引导VCD、CDA视觉迁移和同措辞引导M3ID分别比较四个表达。选择目标为$J=(C+TP)/N$；并列时优先整体准确率，再按UNKNOWN、UNCLEAR、UNSURE、I cannot identify it的顺序选择。表\ref{tab:supp-dev-selection}给出实际选择及其2,424题评估结果。同预算的引导VCD和引导M3ID用于分析表达选择的作用。

开发集选择的IP-VCD在Qwen2.5-VL、Qwen3.5-4B、MiniCPM-V2.6和LLaVA-Mistral上，相对原生VCD的评估$J$分别增加15.68、0.91、3.96和14.23个百分点；Gemma-3下降0.66个百分点。Qwen2.5-VL选择UNSURE后$J=44.02\%$，对应评估集四表达最高观测值为$46.37\%$；其余四模型选择结果的$J$与各自最高观测值相同。这些结果给出了仅使用开发输入确定表达后的实际工作点。

\input{statistics/supplementary/dev_selection_table.tex}
\input{statistics/supplementary/dev_eval_counts_table.tex}

表中U、C、S、I依次表示UNKNOWN、UNCLEAR、UNSURE和I cannot identify it。计数C、W、A分别表示正确回答、错误回答和语义弃权；TP与FP分别为参考支持和参考不支持的弃权。每行满足$C+W+A=2,424$且$TP+FP=A$。

\subsection{DoLa与SID算子设置的同题比较}
\label{sec:supp-author-operators}

每类一张图像组成101题配对集合。DoLa采用作者代码的KL方向选择提前层；SID采用作者贪心分支的未截断对比分数和$\alpha=0.5$\cite{chuang2024dola,huo2025sid}。九个DoLa条件与六个SID条件共包含1,515条回复。表\ref{tab:supp-operator101}比较相同输入上的原计算与作者算子计算，表\ref{tab:supp-same101-methods}给出同一集合上各方法的联合效用。此处统计分母为101；完整评估的统计分母为2,424。

DoLa的总正确数从307/909变为306/909，平均$J$变化为$-0.11$个百分点；SID从203/606变为195/606，平均变化为$-1.32$个百分点。模型间变化具有异质性：Gemma-3的DoLa正确数从48增至53，Phi-3.5-Vision从36减至30。Gemma-3在该集合上的DoLa为53/101，IP-VCD为49/101，IP-M3ID为51/101；其余八模型的IP-VCD均高于DoLa。

\input{statistics/supplementary/operator_table.tex}
\input{statistics/supplementary/same101_comparison_table.tex}

\subsection{CDA视觉迁移的条件权重}
\label{sec:supp-cda-weights}

九个模型各选取12条回复路径，覆盖可用的正确、错误、弃权行为与参考阳性、阴性组合，共108条路径。沿实际回复前缀进行五路教师强制测量，得到803个词元位置。采用CDA的无动量校准式\cite{kim2025cda}，令$H_p,H_c$为先验与视觉上下文分支熵，$H_{0p},H_{0c}$为对应空输入熵，则
\begin{align}
 r_p&=\frac{\max(H_p-H_{0p},0)}{H_{0p}},&
 r_c&=\frac{\max(H_c-H_{0c},0)}{H_{0c}},\\
 w_p&=\frac{r_p^2}{r_p+r_c},&
 w_c&=\frac{r_c^2}{r_p+r_c},\qquad w_a=1-w_p-w_c.
\end{align}
当$r_p+r_c=0$时，连续延拓给出$(w_p,w_c,w_a)=(0,0,1)$。三个权重作用于先验、视觉上下文和弃权条件的logit。

803个位置中，418个使用零和连续延拓，145个具有负弃权权重。表\ref{tab:supp-cda-weights}按回复行为给出逐位置均值。Qwen3-VL的正确、错误和弃权三组均出现负的平均$w_a$；MiniCPM-V2.6和OneVision的所测位置均未出现负$w_a$。Gemma-3错误回答组的平均权重幅度达到$10^5$量级，显示熵比值校准可产生强烈的分支放大。CDA在这些路径上表现为随输入和前缀改变的logit线性组合，权重符号与回答行为的对应关系依赖模型。

\input{statistics/supplementary/cda_behavior_table.tex}

这些均值的统计单位为词元位置，路径按行为分层选取。五路分布计算的权重与公式重算值一致；沿相同前缀的当前argmax在776/803个位置与原回复词元相同，其中Gemma-3有26处差异，Qwen3.5-4B有1处差异。表中权重描述本次前缀测量。

\subsection{共享前缀重放中的噪声与硬件条件}
\label{sec:supp-replay-conditions}

六条回复路径上的九个位置在两类硬件和两种会话顺序下共测量36次。正常视觉输入、提示词元、图像和随机种子保持相同；两类硬件生成的扰动视觉张量不同。3090执行在18/18个配对位置上与原回复词元相同，A100或K100执行为0/18。改变会话顺序后，18/18对argmax保持一致（表\ref{tab:supp-replay}）。

进一步在MiniCPM-V2.6的一条苹果派路径上直接共享同一已保存噪声张量，比较两个词元位置和两种会话顺序。3090的4/4个位置保持原词元，A100的首词元在2/2次执行中恢复一致；第二词元在两种顺序下均由句点变为``with''。共享噪声后，该实例的差异集中在第二词元，表明后续解码的数值路径也参与跨硬件差异。

\input{statistics/supplementary/replay_table.tex}
}{}
\end{document}
